
\documentclass{article}
\usepackage{colortbl}
\usepackage{makecell}
\usepackage{multirow}
\usepackage{supertabular}

\begin{document}

\newcounter{utterance}

\centering \large Interaction Transcript for game `clean\_up', experiment `0\_easy\_7obj\_uk', episode 2 with Qwen3.6{-}35B{-}A3B{-}FP8{-}without{-}reasoning.
\vspace{24pt}

{ \footnotesize  \setcounter{utterance}{1}
\setlength{\tabcolsep}{0pt}
\begin{supertabular}{c@{$\;$}|p{.15\linewidth}@{}p{.15\linewidth}p{.15\linewidth}p{.15\linewidth}p{.15\linewidth}p{.15\linewidth}}
    \# & $\;$A & \multicolumn{4}{c}{Game Master} & $\;\:$B\\
    \hline

    \theutterance \stepcounter{utterance}  
    & & \multicolumn{4}{p{0.6\linewidth}}{
        \cellcolor[rgb]{0.9,0.9,0.9}{
            \makecell[{{p{\linewidth}}}]{
                \texttt{\tiny{[P1$\langle$GM]}}
                \texttt{Я ваш ведучий гри, і ви граєте у спільну гру з наступною сіткою як ігровим полем:} \\
\\ 
\texttt{```} \\
\texttt{ 1234567} \\
\texttt{╔═══════╗} \\
\texttt{║◌T◌◌◌◌◌║ 1} \\
\texttt{║A◌◌◌◌◌D║ 2} \\
\texttt{║◌◌◌◌◌┌─╢ 3} \\
\texttt{╟──┐O◌│◌║ 4} \\
\texttt{║◌P├──┼─╢ 5} \\
\texttt{║◌S│M◌│◌║ 6} \\
\texttt{║◌◌│◌◌│◌║ 7} \\
\texttt{╚══╧══╧═╝} \\
\\ 
\texttt{```} \\
\\ 
\texttt{* Верхній край показує координати x, що зростають праворуч, а правий край — координати y, що зростають донизу.} \\
\texttt{* Наступні об'єкти випадково розміщені на вашій сітці: 'P', 'O', 'T', 'S', 'D', 'A', 'M'.} \\
\\ 
\\ 
\texttt{Інший гравець бачить варіант ігрового поля, де об'єкти розміщені в інших випадкових місцях. Ви не бачите поле іншого гравця, а він не бачить вашого.} \\
\\ 
\\ 
\texttt{**Мета:**} \\
\\ 
\texttt{Обидва гравці повинні переміщати об'єкти на своєму фоні так, щоб однакові об'єкти опинилися в однакових координатах. Ви повинні спілкуватися з іншим гравцем, щоб узгодити спільну цільову конфігурацію.} \\
\\ 
\\ 
\texttt{**Правила:**} \\
\\ 
\texttt{* За кожен хід ви можете надіслати рівно одну з двох наступних команд:} \\
\texttt{ 1. `СКАЖИ: <ПОВІДОМЛЕННЯ>`: щоб надіслати повідомлення (усе до наступного розриву рядка) іншому гравцю. Я передам його вашому партнеру.} \\
\texttt{2. `ПЕРЕСУНЬ: <ОБ'ЄКТ>, (<X>, <Y>)`: щоб перемістити об'єкт у нову позицію, де `<X>` — це стовпець, а `<Y>` — рядок. Я повідомлю вас, чи був ваш хід успішним.} \\
\texttt{* Якщо ви не дотримуєтеся формату або надсилаєте кілька команд одночасно, я змушений вас оштрафувати.} \\
\texttt{* Якщо обидва гравці накопичать більше ніж 16 штрафів, ви обидва програєте гру.} \\
\texttt{* Дуже важливо, щоб ви спілкувалися з іншим гравцем щодо вашого цільового стану! *Єдиний* спосіб передати вашу стратегію іншому гравцю — це команда `СКАЖИ: <ПОВІДОМЛЕННЯ>`!} \\
\\ 
\\ 
\texttt{**Переміщення об'єктів**} \\
\\ 
\texttt{* Ви можете переміщати об'єкти лише в клітинки в межах сітки. Цільова клітинка повинна бути порожньою, тобто містити лише символ '◌'.} \\
\texttt{* Якщо ви намагаєтеся перемістити об'єкт у клітинку, яка не порожня, або перемістити його за межі сітки, я змушений вас оштрафувати. Ви отримаєте ще одну спробу.} \\
\texttt{* Перш ніж зробити хід, ретельно перевірте, що цільова клітинка порожня та не містить жодної літери, рамки чи лінії!} \\
\\ 
\\ 
\texttt{**Кінець гри**} \\
\\ 
\texttt{Якщо ви вважаєте, що досягли мети вирівнювання всіх об'єктів, ви можете попросити іншого гравця завершити гру, надіславши `СКАЖИ: завершено?`. Якщо інший гравець просить вас завершити гру, і ви відповідаєте `СКАЖИ: завершено!`, гра завершиться.} \\
\\ 
\texttt{Обидва гравці перемагають, якщо гра завершується протягом 28 раундів, де один раунд визначається як надсилання обома гравцями по одній дійсній команді.} \\
\\ 
\texttt{**Підрахунок очок:**} \\
\\ 
\texttt{Чим ближче однакові об'єкти на обох ігрових полях, тим більше очок ви отримуєте. Штрафи зменшують ваші очки. Чи зможете ви побити рекорд?} \\
\\ 
\\ 
\texttt{Будь ласка, надішліть повідомлення іншому гравцю, щоб почати гру!} \\
            }
        }
    }
    & \\ \\

    \theutterance \stepcounter{utterance}  
    & \multicolumn{4}{p{0.6\linewidth}}{
        \cellcolor[rgb]{0.9,0.9,0.9}{
            \makecell[{{p{\linewidth}}}]{
                \texttt{\tiny{[P1$\rangle$GM]}}
                \texttt{СКАЖИ: Привіт! Давай домовимось про цільову позицію для всіх фігур. Я пропоную розмістити їх у лівому верхньому куті, щоб уникнути плутанини з рамками.} \\
\\ 
\texttt{Ось моя пропозиція координат (X, Y):} \\
\texttt{T: (1, 1)} \\
\texttt{A: (2, 1)} \\
\texttt{D: (3, 1)} \\
\texttt{O: (1, 2)} \\
\texttt{P: (2, 2)} \\
\texttt{S: (3, 2)} \\
\texttt{M: (1, 3)} \\
\\ 
\texttt{Якщо ти згоден, я почну переміщати свої фігури туди. Напиши "ЗГОДЕН", якщо ця схема тобі підходить, або запропонуй свою.} \\
            }
        }
    }
    & & \\ \\

    \theutterance \stepcounter{utterance}  
    & & & \multicolumn{2}{p{0.3\linewidth}}{
        \cellcolor[rgb]{0.9,0.9,0.9}{
            \makecell[{{p{\linewidth}}}]{
                \texttt{\tiny{[GM$|$GM]}}
                \texttt{СКАЖИ: Привіт! Давай домовимось про цільову позицію для всіх фігур. Я пропоную розмістити їх у лівому верхньому куті, щоб уникнути плутанини з рамками.} \\
\\ 
\texttt{Ось моя пропозиція координат (X, Y):} \\
\texttt{T: (1, 1)} \\
\texttt{A: (2, 1)} \\
\texttt{D: (3, 1)} \\
\texttt{O: (1, 2)} \\
\texttt{P: (2, 2)} \\
\texttt{S: (3, 2)} \\
\texttt{M: (1, 3)} \\
\\ 
\texttt{Якщо ти згоден, я почну переміщати свої фігури туди. Напиши "ЗГОДЕН", якщо ця схема тобі підходить, або запропонуй свою.} \\
            }
        }
    }
    & & \\ \\

    \theutterance \stepcounter{utterance}  
    & & & \multicolumn{2}{p{0.3\linewidth}}{
        \cellcolor[rgb]{0.9,0.9,0.9}{
            \makecell[{{p{\linewidth}}}]{
                \texttt{\tiny{[GM$|$GM]}}
                \texttt{Invalid format: tail is not empty: '} \\
\\ 
\texttt{Ось моя пропозиція координат (X, Y):} \\
\texttt{T: (1, 1)} \\
\texttt{A: (2, 1)} \\
\texttt{D: (3, 1)} \\
\texttt{O: (1, 2)} \\
\texttt{P: (2, 2)} \\
\texttt{S: (3, 2)} \\
\texttt{M: (1, 3)} \\
\\ 
\texttt{Якщо ти згоден, я почну переміщати свої фігури туди. Напиши "ЗГОДЕН", якщо ця схема тобі підходить, або запропонуй свою.} \\
            }
        }
    }
    & & \\ \\

    \theutterance \stepcounter{utterance}  
    & & \multicolumn{4}{p{0.6\linewidth}}{
        \cellcolor[rgb]{0.9,0.9,0.9}{
            \makecell[{{p{\linewidth}}}]{
                \texttt{\tiny{[P1$\langle$GM]}}
                \texttt{Штраф: Ваше повідомлення не повинно містити нічого після команди!} \\
\\ 
\texttt{Переконайтеся, що ваша відповідь містить лише `СКАЖИ: <ПОВІДОМЛЕННЯ>` або `ПЕРЕСУНЬ: <ОБ'ЄКТ>, (<X>, <Y>)` і нічого більше!} \\
\\ 
\texttt{Разом ви накопичили 1 з 16 штрафів.} \\
\texttt{Будь ласка, спробуйте ще раз!} \\
            }
        }
    }
    & \\ \\

    \theutterance \stepcounter{utterance}  
    & \multicolumn{4}{p{0.6\linewidth}}{
        \cellcolor[rgb]{0.9,0.9,0.9}{
            \makecell[{{p{\linewidth}}}]{
                \texttt{\tiny{[P1$\rangle$GM]}}
                \texttt{СКАЖИ: Привіт! Давай домовимось про цільову позицію для всіх фігур. Я пропоную розмістити їх у лівому верхньому куті, щоб уникнути плутанини з рамками. Ось моя пропозиція координат (X, Y): T: (1, 1), A: (2, 1), D: (3, 1), O: (1, 2), P: (2, 2), S: (3, 2), M: (1, 3). Якщо ти згоден, я почну переміщати свої фігури туди. Напиши "ЗГОДЕН", якщо ця схема тобі підходить, або запропонуй свою.} \\
            }
        }
    }
    & & \\ \\

    \theutterance \stepcounter{utterance}  
    & & & \multicolumn{2}{p{0.3\linewidth}}{
        \cellcolor[rgb]{0.9,0.9,0.9}{
            \makecell[{{p{\linewidth}}}]{
                \texttt{\tiny{[GM$|$GM]}}
                \texttt{СКАЖИ: Привіт! Давай домовимось про цільову позицію для всіх фігур. Я пропоную розмістити їх у лівому верхньому куті, щоб уникнути плутанини з рамками. Ось моя пропозиція координат (X, Y): T: (1, 1), A: (2, 1), D: (3, 1), O: (1, 2), P: (2, 2), S: (3, 2), M: (1, 3). Якщо ти згоден, я почну переміщати свої фігури туди. Напиши "ЗГОДЕН", якщо ця схема тобі підходить, або запропонуй свою.} \\
            }
        }
    }
    & & \\ \\

    \theutterance \stepcounter{utterance}  
    & & \multicolumn{4}{p{0.6\linewidth}}{
        \cellcolor[rgb]{0.9,0.9,0.9}{
            \makecell[{{p{\linewidth}}}]{
                \texttt{\tiny{[GM$\rangle$P2]}}
                \texttt{Я ваш ведучий гри, і ви граєте у спільну гру з наступною сіткою як ігровим полем:} \\
\\ 
\texttt{```} \\
\texttt{ 1234567} \\
\texttt{╔═══════╗} \\
\texttt{║◌◌P◌A◌O║ 1} \\
\texttt{║S◌◌◌M◌◌║ 2} \\
\texttt{║◌◌◌◌D┌─╢ 3} \\
\texttt{╟──┐◌◌│◌║ 4} \\
\texttt{║◌◌├──┼─╢ 5} \\
\texttt{║◌◌│◌◌│T║ 6} \\
\texttt{║◌◌│◌◌│◌║ 7} \\
\texttt{╚══╧══╧═╝} \\
\\ 
\texttt{```} \\
\\ 
\texttt{* Верхній край показує координати x, що зростають праворуч, а правий край — координати y, що зростають донизу.} \\
\texttt{* Наступні об'єкти випадково розміщені на вашій сітці: 'P', 'O', 'T', 'S', 'D', 'A', 'M'.} \\
\\ 
\\ 
\texttt{Інший гравець бачить варіант ігрового поля, де об'єкти розміщені в інших випадкових місцях. Ви не бачите поле іншого гравця, а він не бачить вашого.} \\
\\ 
\\ 
\texttt{**Мета:**} \\
\\ 
\texttt{Обидва гравці повинні переміщати об'єкти на своєму фоні так, щоб однакові об'єкти опинилися в однакових координатах. Ви повинні спілкуватися з іншим гравцем, щоб узгодити спільну цільову конфігурацію.} \\
\\ 
\\ 
\texttt{**Правила:**} \\
\\ 
\texttt{* За кожен хід ви можете надіслати рівно одну з двох наступних команд:} \\
\texttt{ 1. `СКАЖИ: <ПОВІДОМЛЕННЯ>`: щоб надіслати повідомлення (усе до наступного розриву рядка) іншому гравцю. Я передам його вашому партнеру.} \\
\texttt{2. `ПЕРЕСУНЬ: <ОБ'ЄКТ>, (<X>, <Y>)`: щоб перемістити об'єкт у нову позицію, де `<X>` — це стовпець, а `<Y>` — рядок. Я повідомлю вас, чи був ваш хід успішним.} \\
\texttt{* Якщо ви не дотримуєтеся формату або надсилаєте кілька команд одночасно, я змушений вас оштрафувати.} \\
\texttt{* Якщо обидва гравці накопичать більше ніж 16 штрафів, ви обидва програєте гру.} \\
\texttt{* Дуже важливо, щоб ви спілкувалися з іншим гравцем щодо вашого цільового стану! *Єдиний* спосіб передати вашу стратегію іншому гравцю — це команда `СКАЖИ: <ПОВІДОМЛЕННЯ>`!} \\
\\ 
\\ 
\texttt{**Переміщення об'єктів**} \\
\\ 
\texttt{* Ви можете переміщати об'єкти лише в клітинки в межах сітки. Цільова клітинка повинна бути порожньою, тобто містити лише символ '◌'.} \\
\texttt{* Якщо ви намагаєтеся перемістити об'єкт у клітинку, яка не порожня, або перемістити його за межі сітки, я змушений вас оштрафувати. Ви отримаєте ще одну спробу.} \\
\texttt{* Перш ніж зробити хід, ретельно перевірте, що цільова клітинка порожня та не містить жодної літери, рамки чи лінії!} \\
\\ 
\\ 
\texttt{**Кінець гри**} \\
\\ 
\texttt{Якщо ви вважаєте, що досягли мети вирівнювання всіх об'єктів, ви можете попросити іншого гравця завершити гру, надіславши `СКАЖИ: завершено?`. Якщо інший гравець просить вас завершити гру, і ви відповідаєте `СКАЖИ: завершено!`, гра завершиться.} \\
\\ 
\texttt{Обидва гравці перемагають, якщо гра завершується протягом 28 раундів, де один раунд визначається як надсилання обома гравцями по одній дійсній команді.} \\
\\ 
\texttt{**Підрахунок очок:**} \\
\\ 
\texttt{Чим ближче однакові об'єкти на обох ігрових полях, тим більше очок ви отримуєте. Штрафи зменшують ваші очки. Чи зможете ви побити рекорд?} \\
\\ 
\texttt{Інший гравець почав гру, надіславши це повідомлення:} \\
\\ 
\texttt{"Привіт! Давай домовимось про цільову позицію для всіх фігур. Я пропоную розмістити їх у лівому верхньому куті, щоб уникнути плутанини з рамками. Ось моя пропозиція координат (X, Y): T: (1, 1), A: (2, 1), D: (3, 1), O: (1, 2), P: (2, 2), S: (3, 2), M: (1, 3). Якщо ти згоден, я почну переміщати свої фігури туди. Напиши "ЗГОДЕН", якщо ця схема тобі підходить, або запропонуй свою."} \\
\\ 
\texttt{Яка ваша перша команда?} \\
            }
        }
    }
    & \\ \\

    \theutterance \stepcounter{utterance}  
    & & & \multicolumn{4}{p{0.6\linewidth}}{
        \cellcolor[rgb]{0.9,0.9,0.9}{
            \makecell[{{p{\linewidth}}}]{
                \texttt{\tiny{[GM$\langle$P2]}}
                \texttt{СКАЖИ: ЗГОДЕН} \\
            }
        }
    }
     \\ \\

    \theutterance \stepcounter{utterance}  
    & & & \multicolumn{2}{p{0.3\linewidth}}{
        \cellcolor[rgb]{0.9,0.9,0.9}{
            \makecell[{{p{\linewidth}}}]{
                \texttt{\tiny{[GM$|$GM]}}
                \texttt{СКАЖИ: ЗГОДЕН} \\
            }
        }
    }
    & & \\ \\

    \theutterance \stepcounter{utterance}  
    & & \multicolumn{4}{p{0.6\linewidth}}{
        \cellcolor[rgb]{0.9,0.9,0.9}{
            \makecell[{{p{\linewidth}}}]{
                \texttt{\tiny{[P1$\langle$GM]}}
                \texttt{Ваше повідомлення передано іншому гравцю.} \\
\\ 
\\ 
\texttt{Зараз ви граєте раунд 1 з максимум 28.} \\
\texttt{Разом ви накопичили 1 з 16 штрафів.} \\
\texttt{Інший гравець надіслав це повідомлення:} \\
\\ 
\\ 
\texttt{"ЗГОДЕН"} \\
\\ 
\\ 
\texttt{Яка ваша наступна команда?} \\
            }
        }
    }
    & \\ \\

    \theutterance \stepcounter{utterance}  
    & \multicolumn{4}{p{0.6\linewidth}}{
        \cellcolor[rgb]{0.9,0.9,0.9}{
            \makecell[{{p{\linewidth}}}]{
                \texttt{\tiny{[P1$\rangle$GM]}}
                \texttt{ПЕРЕСУНЬ: T, (1, 1)} \\
            }
        }
    }
    & & \\ \\

    \theutterance \stepcounter{utterance}  
    & & & \multicolumn{2}{p{0.3\linewidth}}{
        \cellcolor[rgb]{0.9,0.9,0.9}{
            \makecell[{{p{\linewidth}}}]{
                \texttt{\tiny{[GM$|$GM]}}
                \texttt{ПЕРЕСУНЬ: T, (1, 1)} \\
            }
        }
    }
    & & \\ \\

    \theutterance \stepcounter{utterance}  
    & & & \multicolumn{2}{p{0.3\linewidth}}{
        \cellcolor[rgb]{0.9,0.9,0.9}{
            \makecell[{{p{\linewidth}}}]{
                \texttt{\tiny{[GM$|$GM]}}
                \texttt{T moved to (1, 1)} \\
            }
        }
    }
    & & \\ \\

    \theutterance \stepcounter{utterance}  
    & & \multicolumn{4}{p{0.6\linewidth}}{
        \cellcolor[rgb]{0.9,0.9,0.9}{
            \makecell[{{p{\linewidth}}}]{
                \texttt{\tiny{[GM$\rangle$P2]}}
                \texttt{Ваше повідомлення передано іншому гравцю.} \\
\\ 
\\ 
\texttt{Зараз ви граєте раунд 2 з максимум 28.} \\
\texttt{Разом ви накопичили 1 з 16 штрафів.} \\
\texttt{Інший гравець перемістив об'єкт на своїй сітці.} \\
\\ 
\\ 
\texttt{Яка ваша наступна команда?} \\
            }
        }
    }
    & \\ \\

    \theutterance \stepcounter{utterance}  
    & & & \multicolumn{4}{p{0.6\linewidth}}{
        \cellcolor[rgb]{0.9,0.9,0.9}{
            \makecell[{{p{\linewidth}}}]{
                \texttt{\tiny{[GM$\langle$P2]}}
                \texttt{ПЕРЕСУНЬ: P, (1, 1)} \\
            }
        }
    }
     \\ \\

    \theutterance \stepcounter{utterance}  
    & & & \multicolumn{2}{p{0.3\linewidth}}{
        \cellcolor[rgb]{0.9,0.9,0.9}{
            \makecell[{{p{\linewidth}}}]{
                \texttt{\tiny{[GM$|$GM]}}
                \texttt{ПЕРЕСУНЬ: P, (1, 1)} \\
            }
        }
    }
    & & \\ \\

    \theutterance \stepcounter{utterance}  
    & & & \multicolumn{2}{p{0.3\linewidth}}{
        \cellcolor[rgb]{0.9,0.9,0.9}{
            \makecell[{{p{\linewidth}}}]{
                \texttt{\tiny{[GM$|$GM]}}
                \texttt{P moved to (1, 1)} \\
            }
        }
    }
    & & \\ \\

    \theutterance \stepcounter{utterance}  
    & & \multicolumn{4}{p{0.6\linewidth}}{
        \cellcolor[rgb]{0.9,0.9,0.9}{
            \makecell[{{p{\linewidth}}}]{
                \texttt{\tiny{[P1$\langle$GM]}}
                \texttt{Об'єкт 'T' успішно переміщено в (1,1). Ваша оновлена сітка виглядає так:} \\
\texttt{```} \\
\texttt{ 1234567} \\
\texttt{╔═══════╗} \\
\texttt{║T◌◌◌◌◌◌║ 1} \\
\texttt{║A◌◌◌◌◌D║ 2} \\
\texttt{║◌◌◌◌◌┌─╢ 3} \\
\texttt{╟──┐O◌│◌║ 4} \\
\texttt{║◌P├──┼─╢ 5} \\
\texttt{║◌S│M◌│◌║ 6} \\
\texttt{║◌◌│◌◌│◌║ 7} \\
\texttt{╚══╧══╧═╝} \\
\texttt{```} \\
\texttt{Зараз ви граєте раунд 2 з максимум 28.} \\
\texttt{Разом ви накопичили 1 з 16 штрафів.} \\
\texttt{Інший гравець перемістив об'єкт на своїй сітці.} \\
\\ 
\\ 
\texttt{Яка ваша наступна команда?} \\
            }
        }
    }
    & \\ \\

    \theutterance \stepcounter{utterance}  
    & \multicolumn{4}{p{0.6\linewidth}}{
        \cellcolor[rgb]{0.9,0.9,0.9}{
            \makecell[{{p{\linewidth}}}]{
                \texttt{\tiny{[P1$\rangle$GM]}}
                \texttt{ПЕРЕСУНЬ: A, (2, 1)} \\
            }
        }
    }
    & & \\ \\

    \theutterance \stepcounter{utterance}  
    & & & \multicolumn{2}{p{0.3\linewidth}}{
        \cellcolor[rgb]{0.9,0.9,0.9}{
            \makecell[{{p{\linewidth}}}]{
                \texttt{\tiny{[GM$|$GM]}}
                \texttt{ПЕРЕСУНЬ: A, (2, 1)} \\
            }
        }
    }
    & & \\ \\

    \theutterance \stepcounter{utterance}  
    & & & \multicolumn{2}{p{0.3\linewidth}}{
        \cellcolor[rgb]{0.9,0.9,0.9}{
            \makecell[{{p{\linewidth}}}]{
                \texttt{\tiny{[GM$|$GM]}}
                \texttt{A moved to (2, 1)} \\
            }
        }
    }
    & & \\ \\

    \theutterance \stepcounter{utterance}  
    & & \multicolumn{4}{p{0.6\linewidth}}{
        \cellcolor[rgb]{0.9,0.9,0.9}{
            \makecell[{{p{\linewidth}}}]{
                \texttt{\tiny{[GM$\rangle$P2]}}
                \texttt{Об'єкт 'P' успішно переміщено в (1,1). Ваша оновлена сітка виглядає так:} \\
\texttt{```} \\
\texttt{ 1234567} \\
\texttt{╔═══════╗} \\
\texttt{║P◌◌◌A◌O║ 1} \\
\texttt{║S◌◌◌M◌◌║ 2} \\
\texttt{║◌◌◌◌D┌─╢ 3} \\
\texttt{╟──┐◌◌│◌║ 4} \\
\texttt{║◌◌├──┼─╢ 5} \\
\texttt{║◌◌│◌◌│T║ 6} \\
\texttt{║◌◌│◌◌│◌║ 7} \\
\texttt{╚══╧══╧═╝} \\
\texttt{```} \\
\texttt{Зараз ви граєте раунд 3 з максимум 28.} \\
\texttt{Разом ви накопичили 1 з 16 штрафів.} \\
\texttt{Інший гравець перемістив об'єкт на своїй сітці.} \\
\\ 
\\ 
\texttt{Яка ваша наступна команда?} \\
            }
        }
    }
    & \\ \\

    \theutterance \stepcounter{utterance}  
    & & & \multicolumn{4}{p{0.6\linewidth}}{
        \cellcolor[rgb]{0.9,0.9,0.9}{
            \makecell[{{p{\linewidth}}}]{
                \texttt{\tiny{[GM$\langle$P2]}}
                \texttt{ПЕРЕСУНЬ: A, (2, 1)} \\
            }
        }
    }
     \\ \\

    \theutterance \stepcounter{utterance}  
    & & & \multicolumn{2}{p{0.3\linewidth}}{
        \cellcolor[rgb]{0.9,0.9,0.9}{
            \makecell[{{p{\linewidth}}}]{
                \texttt{\tiny{[GM$|$GM]}}
                \texttt{ПЕРЕСУНЬ: A, (2, 1)} \\
            }
        }
    }
    & & \\ \\

    \theutterance \stepcounter{utterance}  
    & & & \multicolumn{2}{p{0.3\linewidth}}{
        \cellcolor[rgb]{0.9,0.9,0.9}{
            \makecell[{{p{\linewidth}}}]{
                \texttt{\tiny{[GM$|$GM]}}
                \texttt{A moved to (2, 1)} \\
            }
        }
    }
    & & \\ \\

    \theutterance \stepcounter{utterance}  
    & & \multicolumn{4}{p{0.6\linewidth}}{
        \cellcolor[rgb]{0.9,0.9,0.9}{
            \makecell[{{p{\linewidth}}}]{
                \texttt{\tiny{[P1$\langle$GM]}}
                \texttt{Об'єкт 'A' успішно переміщено в (2,1). Ваша оновлена сітка виглядає так:} \\
\texttt{```} \\
\texttt{ 1234567} \\
\texttt{╔═══════╗} \\
\texttt{║TA◌◌◌◌◌║ 1} \\
\texttt{║◌◌◌◌◌◌D║ 2} \\
\texttt{║◌◌◌◌◌┌─╢ 3} \\
\texttt{╟──┐O◌│◌║ 4} \\
\texttt{║◌P├──┼─╢ 5} \\
\texttt{║◌S│M◌│◌║ 6} \\
\texttt{║◌◌│◌◌│◌║ 7} \\
\texttt{╚══╧══╧═╝} \\
\texttt{```} \\
\texttt{Зараз ви граєте раунд 3 з максимум 28.} \\
\texttt{Разом ви накопичили 1 з 16 штрафів.} \\
\texttt{Інший гравець перемістив об'єкт на своїй сітці.} \\
\\ 
\\ 
\texttt{Яка ваша наступна команда?} \\
            }
        }
    }
    & \\ \\

    \theutterance \stepcounter{utterance}  
    & \multicolumn{4}{p{0.6\linewidth}}{
        \cellcolor[rgb]{0.9,0.9,0.9}{
            \makecell[{{p{\linewidth}}}]{
                \texttt{\tiny{[P1$\rangle$GM]}}
                \texttt{ПЕРЕСУНЬ: D, (3, 1)} \\
            }
        }
    }
    & & \\ \\

    \theutterance \stepcounter{utterance}  
    & & & \multicolumn{2}{p{0.3\linewidth}}{
        \cellcolor[rgb]{0.9,0.9,0.9}{
            \makecell[{{p{\linewidth}}}]{
                \texttt{\tiny{[GM$|$GM]}}
                \texttt{ПЕРЕСУНЬ: D, (3, 1)} \\
            }
        }
    }
    & & \\ \\

    \theutterance \stepcounter{utterance}  
    & & & \multicolumn{2}{p{0.3\linewidth}}{
        \cellcolor[rgb]{0.9,0.9,0.9}{
            \makecell[{{p{\linewidth}}}]{
                \texttt{\tiny{[GM$|$GM]}}
                \texttt{D moved to (3, 1)} \\
            }
        }
    }
    & & \\ \\

    \theutterance \stepcounter{utterance}  
    & & \multicolumn{4}{p{0.6\linewidth}}{
        \cellcolor[rgb]{0.9,0.9,0.9}{
            \makecell[{{p{\linewidth}}}]{
                \texttt{\tiny{[GM$\rangle$P2]}}
                \texttt{Об'єкт 'A' успішно переміщено в (2,1). Ваша оновлена сітка виглядає так:} \\
\texttt{```} \\
\texttt{ 1234567} \\
\texttt{╔═══════╗} \\
\texttt{║PA◌◌◌◌O║ 1} \\
\texttt{║S◌◌◌M◌◌║ 2} \\
\texttt{║◌◌◌◌D┌─╢ 3} \\
\texttt{╟──┐◌◌│◌║ 4} \\
\texttt{║◌◌├──┼─╢ 5} \\
\texttt{║◌◌│◌◌│T║ 6} \\
\texttt{║◌◌│◌◌│◌║ 7} \\
\texttt{╚══╧══╧═╝} \\
\texttt{```} \\
\texttt{Зараз ви граєте раунд 4 з максимум 28.} \\
\texttt{Разом ви накопичили 1 з 16 штрафів.} \\
\texttt{Інший гравець перемістив об'єкт на своїй сітці.} \\
\\ 
\\ 
\texttt{Яка ваша наступна команда?} \\
            }
        }
    }
    & \\ \\

    \theutterance \stepcounter{utterance}  
    & & & \multicolumn{4}{p{0.6\linewidth}}{
        \cellcolor[rgb]{0.9,0.9,0.9}{
            \makecell[{{p{\linewidth}}}]{
                \texttt{\tiny{[GM$\langle$P2]}}
                \texttt{ПЕРЕСУНЬ: D, (3, 1)} \\
            }
        }
    }
     \\ \\

    \theutterance \stepcounter{utterance}  
    & & & \multicolumn{2}{p{0.3\linewidth}}{
        \cellcolor[rgb]{0.9,0.9,0.9}{
            \makecell[{{p{\linewidth}}}]{
                \texttt{\tiny{[GM$|$GM]}}
                \texttt{ПЕРЕСУНЬ: D, (3, 1)} \\
            }
        }
    }
    & & \\ \\

    \theutterance \stepcounter{utterance}  
    & & & \multicolumn{2}{p{0.3\linewidth}}{
        \cellcolor[rgb]{0.9,0.9,0.9}{
            \makecell[{{p{\linewidth}}}]{
                \texttt{\tiny{[GM$|$GM]}}
                \texttt{D moved to (3, 1)} \\
            }
        }
    }
    & & \\ \\

    \theutterance \stepcounter{utterance}  
    & & \multicolumn{4}{p{0.6\linewidth}}{
        \cellcolor[rgb]{0.9,0.9,0.9}{
            \makecell[{{p{\linewidth}}}]{
                \texttt{\tiny{[P1$\langle$GM]}}
                \texttt{Об'єкт 'D' успішно переміщено в (3,1). Ваша оновлена сітка виглядає так:} \\
\texttt{```} \\
\texttt{ 1234567} \\
\texttt{╔═══════╗} \\
\texttt{║TAD◌◌◌◌║ 1} \\
\texttt{║◌◌◌◌◌◌◌║ 2} \\
\texttt{║◌◌◌◌◌┌─╢ 3} \\
\texttt{╟──┐O◌│◌║ 4} \\
\texttt{║◌P├──┼─╢ 5} \\
\texttt{║◌S│M◌│◌║ 6} \\
\texttt{║◌◌│◌◌│◌║ 7} \\
\texttt{╚══╧══╧═╝} \\
\texttt{```} \\
\texttt{Зараз ви граєте раунд 4 з максимум 28.} \\
\texttt{Разом ви накопичили 1 з 16 штрафів.} \\
\texttt{Інший гравець перемістив об'єкт на своїй сітці.} \\
\\ 
\\ 
\texttt{Яка ваша наступна команда?} \\
            }
        }
    }
    & \\ \\

    \theutterance \stepcounter{utterance}  
    & \multicolumn{4}{p{0.6\linewidth}}{
        \cellcolor[rgb]{0.9,0.9,0.9}{
            \makecell[{{p{\linewidth}}}]{
                \texttt{\tiny{[P1$\rangle$GM]}}
                \texttt{ПЕРЕСУНЬ: O, (1, 2)} \\
            }
        }
    }
    & & \\ \\

    \theutterance \stepcounter{utterance}  
    & & & \multicolumn{2}{p{0.3\linewidth}}{
        \cellcolor[rgb]{0.9,0.9,0.9}{
            \makecell[{{p{\linewidth}}}]{
                \texttt{\tiny{[GM$|$GM]}}
                \texttt{ПЕРЕСУНЬ: O, (1, 2)} \\
            }
        }
    }
    & & \\ \\

    \theutterance \stepcounter{utterance}  
    & & & \multicolumn{2}{p{0.3\linewidth}}{
        \cellcolor[rgb]{0.9,0.9,0.9}{
            \makecell[{{p{\linewidth}}}]{
                \texttt{\tiny{[GM$|$GM]}}
                \texttt{O moved to (1, 2)} \\
            }
        }
    }
    & & \\ \\

    \theutterance \stepcounter{utterance}  
    & & \multicolumn{4}{p{0.6\linewidth}}{
        \cellcolor[rgb]{0.9,0.9,0.9}{
            \makecell[{{p{\linewidth}}}]{
                \texttt{\tiny{[GM$\rangle$P2]}}
                \texttt{Об'єкт 'D' успішно переміщено в (3,1). Ваша оновлена сітка виглядає так:} \\
\texttt{```} \\
\texttt{ 1234567} \\
\texttt{╔═══════╗} \\
\texttt{║PAD◌◌◌O║ 1} \\
\texttt{║S◌◌◌M◌◌║ 2} \\
\texttt{║◌◌◌◌◌┌─╢ 3} \\
\texttt{╟──┐◌◌│◌║ 4} \\
\texttt{║◌◌├──┼─╢ 5} \\
\texttt{║◌◌│◌◌│T║ 6} \\
\texttt{║◌◌│◌◌│◌║ 7} \\
\texttt{╚══╧══╧═╝} \\
\texttt{```} \\
\texttt{Зараз ви граєте раунд 5 з максимум 28.} \\
\texttt{Разом ви накопичили 1 з 16 штрафів.} \\
\texttt{Інший гравець перемістив об'єкт на своїй сітці.} \\
\\ 
\\ 
\texttt{Яка ваша наступна команда?} \\
            }
        }
    }
    & \\ \\

    \theutterance \stepcounter{utterance}  
    & & & \multicolumn{4}{p{0.6\linewidth}}{
        \cellcolor[rgb]{0.9,0.9,0.9}{
            \makecell[{{p{\linewidth}}}]{
                \texttt{\tiny{[GM$\langle$P2]}}
                \texttt{ПЕРЕСУНЬ: O, (1, 2)} \\
            }
        }
    }
     \\ \\

    \theutterance \stepcounter{utterance}  
    & & & \multicolumn{2}{p{0.3\linewidth}}{
        \cellcolor[rgb]{0.9,0.9,0.9}{
            \makecell[{{p{\linewidth}}}]{
                \texttt{\tiny{[GM$|$GM]}}
                \texttt{ПЕРЕСУНЬ: O, (1, 2)} \\
            }
        }
    }
    & & \\ \\

    \theutterance \stepcounter{utterance}  
    & & & \multicolumn{2}{p{0.3\linewidth}}{
        \cellcolor[rgb]{0.9,0.9,0.9}{
            \makecell[{{p{\linewidth}}}]{
                \texttt{\tiny{[GM$|$GM]}}
                \texttt{Штраф: (1,2) не порожня, а містить 'S'.} \\
\\ 
\texttt{Разом ви накопичили 2 з 16 штрафів.} \\
\texttt{Будь ласка, спробуйте ще раз!} \\
            }
        }
    }
    & & \\ \\

    \theutterance \stepcounter{utterance}  
    & & \multicolumn{4}{p{0.6\linewidth}}{
        \cellcolor[rgb]{0.9,0.9,0.9}{
            \makecell[{{p{\linewidth}}}]{
                \texttt{\tiny{[GM$\rangle$P2]}}
                \texttt{Штраф: (1,2) не порожня, а містить 'S'.} \\
\\ 
\texttt{Разом ви накопичили 2 з 16 штрафів.} \\
\texttt{Будь ласка, спробуйте ще раз!} \\
            }
        }
    }
    & \\ \\

    \theutterance \stepcounter{utterance}  
    & & & \multicolumn{4}{p{0.6\linewidth}}{
        \cellcolor[rgb]{0.9,0.9,0.9}{
            \makecell[{{p{\linewidth}}}]{
                \texttt{\tiny{[GM$\langle$P2]}}
                \texttt{ПЕРЕСУНЬ: O, (1, 3)} \\
            }
        }
    }
     \\ \\

    \theutterance \stepcounter{utterance}  
    & & & \multicolumn{2}{p{0.3\linewidth}}{
        \cellcolor[rgb]{0.9,0.9,0.9}{
            \makecell[{{p{\linewidth}}}]{
                \texttt{\tiny{[GM$|$GM]}}
                \texttt{ПЕРЕСУНЬ: O, (1, 3)} \\
            }
        }
    }
    & & \\ \\

    \theutterance \stepcounter{utterance}  
    & & & \multicolumn{2}{p{0.3\linewidth}}{
        \cellcolor[rgb]{0.9,0.9,0.9}{
            \makecell[{{p{\linewidth}}}]{
                \texttt{\tiny{[GM$|$GM]}}
                \texttt{O moved to (1, 3)} \\
            }
        }
    }
    & & \\ \\

    \theutterance \stepcounter{utterance}  
    & & \multicolumn{4}{p{0.6\linewidth}}{
        \cellcolor[rgb]{0.9,0.9,0.9}{
            \makecell[{{p{\linewidth}}}]{
                \texttt{\tiny{[P1$\langle$GM]}}
                \texttt{Об'єкт 'O' успішно переміщено в (1,2). Ваша оновлена сітка виглядає так:} \\
\texttt{```} \\
\texttt{ 1234567} \\
\texttt{╔═══════╗} \\
\texttt{║TAD◌◌◌◌║ 1} \\
\texttt{║O◌◌◌◌◌◌║ 2} \\
\texttt{║◌◌◌◌◌┌─╢ 3} \\
\texttt{╟──┐◌◌│◌║ 4} \\
\texttt{║◌P├──┼─╢ 5} \\
\texttt{║◌S│M◌│◌║ 6} \\
\texttt{║◌◌│◌◌│◌║ 7} \\
\texttt{╚══╧══╧═╝} \\
\texttt{```} \\
\texttt{Зараз ви граєте раунд 5 з максимум 28.} \\
\texttt{Разом ви накопичили 2 з 16 штрафів.} \\
\texttt{Інший гравець перемістив об'єкт на своїй сітці.} \\
\\ 
\\ 
\texttt{Яка ваша наступна команда?} \\
            }
        }
    }
    & \\ \\

    \theutterance \stepcounter{utterance}  
    & \multicolumn{4}{p{0.6\linewidth}}{
        \cellcolor[rgb]{0.9,0.9,0.9}{
            \makecell[{{p{\linewidth}}}]{
                \texttt{\tiny{[P1$\rangle$GM]}}
                \texttt{ПЕРЕСУНЬ: P, (2, 2)} \\
            }
        }
    }
    & & \\ \\

    \theutterance \stepcounter{utterance}  
    & & & \multicolumn{2}{p{0.3\linewidth}}{
        \cellcolor[rgb]{0.9,0.9,0.9}{
            \makecell[{{p{\linewidth}}}]{
                \texttt{\tiny{[GM$|$GM]}}
                \texttt{ПЕРЕСУНЬ: P, (2, 2)} \\
            }
        }
    }
    & & \\ \\

    \theutterance \stepcounter{utterance}  
    & & & \multicolumn{2}{p{0.3\linewidth}}{
        \cellcolor[rgb]{0.9,0.9,0.9}{
            \makecell[{{p{\linewidth}}}]{
                \texttt{\tiny{[GM$|$GM]}}
                \texttt{P moved to (2, 2)} \\
            }
        }
    }
    & & \\ \\

    \theutterance \stepcounter{utterance}  
    & & \multicolumn{4}{p{0.6\linewidth}}{
        \cellcolor[rgb]{0.9,0.9,0.9}{
            \makecell[{{p{\linewidth}}}]{
                \texttt{\tiny{[GM$\rangle$P2]}}
                \texttt{Об'єкт 'O' успішно переміщено в (1,3). Ваша оновлена сітка виглядає так:} \\
\texttt{```} \\
\texttt{ 1234567} \\
\texttt{╔═══════╗} \\
\texttt{║PAD◌◌◌◌║ 1} \\
\texttt{║S◌◌◌M◌◌║ 2} \\
\texttt{║O◌◌◌◌┌─╢ 3} \\
\texttt{╟──┐◌◌│◌║ 4} \\
\texttt{║◌◌├──┼─╢ 5} \\
\texttt{║◌◌│◌◌│T║ 6} \\
\texttt{║◌◌│◌◌│◌║ 7} \\
\texttt{╚══╧══╧═╝} \\
\texttt{```} \\
\texttt{Зараз ви граєте раунд 6 з максимум 28.} \\
\texttt{Разом ви накопичили 2 з 16 штрафів.} \\
\texttt{Інший гравець перемістив об'єкт на своїй сітці.} \\
\\ 
\\ 
\texttt{Яка ваша наступна команда?} \\
            }
        }
    }
    & \\ \\

    \theutterance \stepcounter{utterance}  
    & & & \multicolumn{4}{p{0.6\linewidth}}{
        \cellcolor[rgb]{0.9,0.9,0.9}{
            \makecell[{{p{\linewidth}}}]{
                \texttt{\tiny{[GM$\langle$P2]}}
                \texttt{ПЕРЕСУНЬ: S, (3, 2)} \\
            }
        }
    }
     \\ \\

    \theutterance \stepcounter{utterance}  
    & & & \multicolumn{2}{p{0.3\linewidth}}{
        \cellcolor[rgb]{0.9,0.9,0.9}{
            \makecell[{{p{\linewidth}}}]{
                \texttt{\tiny{[GM$|$GM]}}
                \texttt{ПЕРЕСУНЬ: S, (3, 2)} \\
            }
        }
    }
    & & \\ \\

    \theutterance \stepcounter{utterance}  
    & & & \multicolumn{2}{p{0.3\linewidth}}{
        \cellcolor[rgb]{0.9,0.9,0.9}{
            \makecell[{{p{\linewidth}}}]{
                \texttt{\tiny{[GM$|$GM]}}
                \texttt{S moved to (3, 2)} \\
            }
        }
    }
    & & \\ \\

    \theutterance \stepcounter{utterance}  
    & & \multicolumn{4}{p{0.6\linewidth}}{
        \cellcolor[rgb]{0.9,0.9,0.9}{
            \makecell[{{p{\linewidth}}}]{
                \texttt{\tiny{[P1$\langle$GM]}}
                \texttt{Об'єкт 'P' успішно переміщено в (2,2). Ваша оновлена сітка виглядає так:} \\
\texttt{```} \\
\texttt{ 1234567} \\
\texttt{╔═══════╗} \\
\texttt{║TAD◌◌◌◌║ 1} \\
\texttt{║OP◌◌◌◌◌║ 2} \\
\texttt{║◌◌◌◌◌┌─╢ 3} \\
\texttt{╟──┐◌◌│◌║ 4} \\
\texttt{║◌◌├──┼─╢ 5} \\
\texttt{║◌S│M◌│◌║ 6} \\
\texttt{║◌◌│◌◌│◌║ 7} \\
\texttt{╚══╧══╧═╝} \\
\texttt{```} \\
\texttt{Зараз ви граєте раунд 6 з максимум 28.} \\
\texttt{Разом ви накопичили 2 з 16 штрафів.} \\
\texttt{Інший гравець перемістив об'єкт на своїй сітці.} \\
\\ 
\\ 
\texttt{Яка ваша наступна команда?} \\
            }
        }
    }
    & \\ \\

    \theutterance \stepcounter{utterance}  
    & \multicolumn{4}{p{0.6\linewidth}}{
        \cellcolor[rgb]{0.9,0.9,0.9}{
            \makecell[{{p{\linewidth}}}]{
                \texttt{\tiny{[P1$\rangle$GM]}}
                \texttt{ПЕРЕСУНЬ: S, (3, 2)} \\
            }
        }
    }
    & & \\ \\

    \theutterance \stepcounter{utterance}  
    & & & \multicolumn{2}{p{0.3\linewidth}}{
        \cellcolor[rgb]{0.9,0.9,0.9}{
            \makecell[{{p{\linewidth}}}]{
                \texttt{\tiny{[GM$|$GM]}}
                \texttt{ПЕРЕСУНЬ: S, (3, 2)} \\
            }
        }
    }
    & & \\ \\

    \theutterance \stepcounter{utterance}  
    & & & \multicolumn{2}{p{0.3\linewidth}}{
        \cellcolor[rgb]{0.9,0.9,0.9}{
            \makecell[{{p{\linewidth}}}]{
                \texttt{\tiny{[GM$|$GM]}}
                \texttt{S moved to (3, 2)} \\
            }
        }
    }
    & & \\ \\

    \theutterance \stepcounter{utterance}  
    & & \multicolumn{4}{p{0.6\linewidth}}{
        \cellcolor[rgb]{0.9,0.9,0.9}{
            \makecell[{{p{\linewidth}}}]{
                \texttt{\tiny{[GM$\rangle$P2]}}
                \texttt{Об'єкт 'S' успішно переміщено в (3,2). Ваша оновлена сітка виглядає так:} \\
\texttt{```} \\
\texttt{ 1234567} \\
\texttt{╔═══════╗} \\
\texttt{║PAD◌◌◌◌║ 1} \\
\texttt{║◌◌S◌M◌◌║ 2} \\
\texttt{║O◌◌◌◌┌─╢ 3} \\
\texttt{╟──┐◌◌│◌║ 4} \\
\texttt{║◌◌├──┼─╢ 5} \\
\texttt{║◌◌│◌◌│T║ 6} \\
\texttt{║◌◌│◌◌│◌║ 7} \\
\texttt{╚══╧══╧═╝} \\
\texttt{```} \\
\texttt{Зараз ви граєте раунд 7 з максимум 28.} \\
\texttt{Разом ви накопичили 2 з 16 штрафів.} \\
\texttt{Інший гравець перемістив об'єкт на своїй сітці.} \\
\\ 
\\ 
\texttt{Яка ваша наступна команда?} \\
            }
        }
    }
    & \\ \\

    \theutterance \stepcounter{utterance}  
    & & & \multicolumn{4}{p{0.6\linewidth}}{
        \cellcolor[rgb]{0.9,0.9,0.9}{
            \makecell[{{p{\linewidth}}}]{
                \texttt{\tiny{[GM$\langle$P2]}}
                \texttt{ПЕРЕСУНЬ: M, (1, 3)} \\
            }
        }
    }
     \\ \\

    \theutterance \stepcounter{utterance}  
    & & & \multicolumn{2}{p{0.3\linewidth}}{
        \cellcolor[rgb]{0.9,0.9,0.9}{
            \makecell[{{p{\linewidth}}}]{
                \texttt{\tiny{[GM$|$GM]}}
                \texttt{ПЕРЕСУНЬ: M, (1, 3)} \\
            }
        }
    }
    & & \\ \\

    \theutterance \stepcounter{utterance}  
    & & & \multicolumn{2}{p{0.3\linewidth}}{
        \cellcolor[rgb]{0.9,0.9,0.9}{
            \makecell[{{p{\linewidth}}}]{
                \texttt{\tiny{[GM$|$GM]}}
                \texttt{Штраф: (1,3) не порожня, а містить 'O'.} \\
\\ 
\texttt{Разом ви накопичили 3 з 16 штрафів.} \\
\texttt{Будь ласка, спробуйте ще раз!} \\
            }
        }
    }
    & & \\ \\

    \theutterance \stepcounter{utterance}  
    & & \multicolumn{4}{p{0.6\linewidth}}{
        \cellcolor[rgb]{0.9,0.9,0.9}{
            \makecell[{{p{\linewidth}}}]{
                \texttt{\tiny{[GM$\rangle$P2]}}
                \texttt{Штраф: (1,3) не порожня, а містить 'O'.} \\
\\ 
\texttt{Разом ви накопичили 3 з 16 штрафів.} \\
\texttt{Будь ласка, спробуйте ще раз!} \\
            }
        }
    }
    & \\ \\

    \theutterance \stepcounter{utterance}  
    & & & \multicolumn{4}{p{0.6\linewidth}}{
        \cellcolor[rgb]{0.9,0.9,0.9}{
            \makecell[{{p{\linewidth}}}]{
                \texttt{\tiny{[GM$\langle$P2]}}
                \texttt{ПЕРЕСУНЬ: M, (2, 3)} \\
            }
        }
    }
     \\ \\

    \theutterance \stepcounter{utterance}  
    & & & \multicolumn{2}{p{0.3\linewidth}}{
        \cellcolor[rgb]{0.9,0.9,0.9}{
            \makecell[{{p{\linewidth}}}]{
                \texttt{\tiny{[GM$|$GM]}}
                \texttt{ПЕРЕСУНЬ: M, (2, 3)} \\
            }
        }
    }
    & & \\ \\

    \theutterance \stepcounter{utterance}  
    & & & \multicolumn{2}{p{0.3\linewidth}}{
        \cellcolor[rgb]{0.9,0.9,0.9}{
            \makecell[{{p{\linewidth}}}]{
                \texttt{\tiny{[GM$|$GM]}}
                \texttt{M moved to (2, 3)} \\
            }
        }
    }
    & & \\ \\

    \theutterance \stepcounter{utterance}  
    & & \multicolumn{4}{p{0.6\linewidth}}{
        \cellcolor[rgb]{0.9,0.9,0.9}{
            \makecell[{{p{\linewidth}}}]{
                \texttt{\tiny{[P1$\langle$GM]}}
                \texttt{Об'єкт 'S' успішно переміщено в (3,2). Ваша оновлена сітка виглядає так:} \\
\texttt{```} \\
\texttt{ 1234567} \\
\texttt{╔═══════╗} \\
\texttt{║TAD◌◌◌◌║ 1} \\
\texttt{║OPS◌◌◌◌║ 2} \\
\texttt{║◌◌◌◌◌┌─╢ 3} \\
\texttt{╟──┐◌◌│◌║ 4} \\
\texttt{║◌◌├──┼─╢ 5} \\
\texttt{║◌◌│M◌│◌║ 6} \\
\texttt{║◌◌│◌◌│◌║ 7} \\
\texttt{╚══╧══╧═╝} \\
\texttt{```} \\
\texttt{Зараз ви граєте раунд 7 з максимум 28.} \\
\texttt{Разом ви накопичили 3 з 16 штрафів.} \\
\texttt{Інший гравець перемістив об'єкт на своїй сітці.} \\
\\ 
\\ 
\texttt{Яка ваша наступна команда?} \\
            }
        }
    }
    & \\ \\

    \theutterance \stepcounter{utterance}  
    & \multicolumn{4}{p{0.6\linewidth}}{
        \cellcolor[rgb]{0.9,0.9,0.9}{
            \makecell[{{p{\linewidth}}}]{
                \texttt{\tiny{[P1$\rangle$GM]}}
                \texttt{ПЕРЕСУНЬ: M, (1, 3)} \\
            }
        }
    }
    & & \\ \\

    \theutterance \stepcounter{utterance}  
    & & & \multicolumn{2}{p{0.3\linewidth}}{
        \cellcolor[rgb]{0.9,0.9,0.9}{
            \makecell[{{p{\linewidth}}}]{
                \texttt{\tiny{[GM$|$GM]}}
                \texttt{ПЕРЕСУНЬ: M, (1, 3)} \\
            }
        }
    }
    & & \\ \\

    \theutterance \stepcounter{utterance}  
    & & & \multicolumn{2}{p{0.3\linewidth}}{
        \cellcolor[rgb]{0.9,0.9,0.9}{
            \makecell[{{p{\linewidth}}}]{
                \texttt{\tiny{[GM$|$GM]}}
                \texttt{M moved to (1, 3)} \\
            }
        }
    }
    & & \\ \\

    \theutterance \stepcounter{utterance}  
    & & \multicolumn{4}{p{0.6\linewidth}}{
        \cellcolor[rgb]{0.9,0.9,0.9}{
            \makecell[{{p{\linewidth}}}]{
                \texttt{\tiny{[GM$\rangle$P2]}}
                \texttt{Об'єкт 'M' успішно переміщено в (2,3). Ваша оновлена сітка виглядає так:} \\
\texttt{```} \\
\texttt{ 1234567} \\
\texttt{╔═══════╗} \\
\texttt{║PAD◌◌◌◌║ 1} \\
\texttt{║◌◌S◌◌◌◌║ 2} \\
\texttt{║OM◌◌◌┌─╢ 3} \\
\texttt{╟──┐◌◌│◌║ 4} \\
\texttt{║◌◌├──┼─╢ 5} \\
\texttt{║◌◌│◌◌│T║ 6} \\
\texttt{║◌◌│◌◌│◌║ 7} \\
\texttt{╚══╧══╧═╝} \\
\texttt{```} \\
\texttt{Зараз ви граєте раунд 8 з максимум 28.} \\
\texttt{Разом ви накопичили 3 з 16 штрафів.} \\
\texttt{Інший гравець перемістив об'єкт на своїй сітці.} \\
\\ 
\\ 
\texttt{Яка ваша наступна команда?} \\
            }
        }
    }
    & \\ \\

    \theutterance \stepcounter{utterance}  
    & & & \multicolumn{4}{p{0.6\linewidth}}{
        \cellcolor[rgb]{0.9,0.9,0.9}{
            \makecell[{{p{\linewidth}}}]{
                \texttt{\tiny{[GM$\langle$P2]}}
                \texttt{ПЕРЕСУНЬ: T, (1, 1)} \\
            }
        }
    }
     \\ \\

    \theutterance \stepcounter{utterance}  
    & & & \multicolumn{2}{p{0.3\linewidth}}{
        \cellcolor[rgb]{0.9,0.9,0.9}{
            \makecell[{{p{\linewidth}}}]{
                \texttt{\tiny{[GM$|$GM]}}
                \texttt{ПЕРЕСУНЬ: T, (1, 1)} \\
            }
        }
    }
    & & \\ \\

    \theutterance \stepcounter{utterance}  
    & & & \multicolumn{2}{p{0.3\linewidth}}{
        \cellcolor[rgb]{0.9,0.9,0.9}{
            \makecell[{{p{\linewidth}}}]{
                \texttt{\tiny{[GM$|$GM]}}
                \texttt{Штраф: (1,1) не порожня, а містить 'P'.} \\
\\ 
\texttt{Разом ви накопичили 4 з 16 штрафів.} \\
\texttt{Будь ласка, спробуйте ще раз!} \\
            }
        }
    }
    & & \\ \\

    \theutterance \stepcounter{utterance}  
    & & \multicolumn{4}{p{0.6\linewidth}}{
        \cellcolor[rgb]{0.9,0.9,0.9}{
            \makecell[{{p{\linewidth}}}]{
                \texttt{\tiny{[GM$\rangle$P2]}}
                \texttt{Штраф: (1,1) не порожня, а містить 'P'.} \\
\\ 
\texttt{Разом ви накопичили 4 з 16 штрафів.} \\
\texttt{Будь ласка, спробуйте ще раз!} \\
            }
        }
    }
    & \\ \\

    \theutterance \stepcounter{utterance}  
    & & & \multicolumn{4}{p{0.6\linewidth}}{
        \cellcolor[rgb]{0.9,0.9,0.9}{
            \makecell[{{p{\linewidth}}}]{
                \texttt{\tiny{[GM$\langle$P2]}}
                \texttt{ПЕРЕСУНЬ: T, (2, 1)} \\
            }
        }
    }
     \\ \\

    \theutterance \stepcounter{utterance}  
    & & & \multicolumn{2}{p{0.3\linewidth}}{
        \cellcolor[rgb]{0.9,0.9,0.9}{
            \makecell[{{p{\linewidth}}}]{
                \texttt{\tiny{[GM$|$GM]}}
                \texttt{ПЕРЕСУНЬ: T, (2, 1)} \\
            }
        }
    }
    & & \\ \\

    \theutterance \stepcounter{utterance}  
    & & & \multicolumn{2}{p{0.3\linewidth}}{
        \cellcolor[rgb]{0.9,0.9,0.9}{
            \makecell[{{p{\linewidth}}}]{
                \texttt{\tiny{[GM$|$GM]}}
                \texttt{Штраф: (2,1) не порожня, а містить 'A'.} \\
\\ 
\texttt{Разом ви накопичили 5 з 16 штрафів.} \\
\texttt{Будь ласка, спробуйте ще раз!} \\
            }
        }
    }
    & & \\ \\

    \theutterance \stepcounter{utterance}  
    & & \multicolumn{4}{p{0.6\linewidth}}{
        \cellcolor[rgb]{0.9,0.9,0.9}{
            \makecell[{{p{\linewidth}}}]{
                \texttt{\tiny{[GM$\rangle$P2]}}
                \texttt{Штраф: (2,1) не порожня, а містить 'A'.} \\
\\ 
\texttt{Разом ви накопичили 5 з 16 штрафів.} \\
\texttt{Будь ласка, спробуйте ще раз!} \\
            }
        }
    }
    & \\ \\

    \theutterance \stepcounter{utterance}  
    & & & \multicolumn{4}{p{0.6\linewidth}}{
        \cellcolor[rgb]{0.9,0.9,0.9}{
            \makecell[{{p{\linewidth}}}]{
                \texttt{\tiny{[GM$\langle$P2]}}
                \texttt{ПЕРЕСУНЬ: T, (3, 1)} \\
            }
        }
    }
     \\ \\

    \theutterance \stepcounter{utterance}  
    & & & \multicolumn{2}{p{0.3\linewidth}}{
        \cellcolor[rgb]{0.9,0.9,0.9}{
            \makecell[{{p{\linewidth}}}]{
                \texttt{\tiny{[GM$|$GM]}}
                \texttt{ПЕРЕСУНЬ: T, (3, 1)} \\
            }
        }
    }
    & & \\ \\

    \theutterance \stepcounter{utterance}  
    & & & \multicolumn{2}{p{0.3\linewidth}}{
        \cellcolor[rgb]{0.9,0.9,0.9}{
            \makecell[{{p{\linewidth}}}]{
                \texttt{\tiny{[GM$|$GM]}}
                \texttt{Штраф: (3,1) не порожня, а містить 'D'.} \\
\\ 
\texttt{Разом ви накопичили 6 з 16 штрафів.} \\
\texttt{Будь ласка, спробуйте ще раз!} \\
            }
        }
    }
    & & \\ \\

    \theutterance \stepcounter{utterance}  
    & & \multicolumn{4}{p{0.6\linewidth}}{
        \cellcolor[rgb]{0.9,0.9,0.9}{
            \makecell[{{p{\linewidth}}}]{
                \texttt{\tiny{[GM$\rangle$P2]}}
                \texttt{Штраф: (3,1) не порожня, а містить 'D'.} \\
\\ 
\texttt{Разом ви накопичили 6 з 16 штрафів.} \\
\texttt{Будь ласка, спробуйте ще раз!} \\
            }
        }
    }
    & \\ \\

    \theutterance \stepcounter{utterance}  
    & & & \multicolumn{4}{p{0.6\linewidth}}{
        \cellcolor[rgb]{0.9,0.9,0.9}{
            \makecell[{{p{\linewidth}}}]{
                \texttt{\tiny{[GM$\langle$P2]}}
                \texttt{ПЕРЕСУНЬ: T, (1, 2)} \\
            }
        }
    }
     \\ \\

    \theutterance \stepcounter{utterance}  
    & & & \multicolumn{2}{p{0.3\linewidth}}{
        \cellcolor[rgb]{0.9,0.9,0.9}{
            \makecell[{{p{\linewidth}}}]{
                \texttt{\tiny{[GM$|$GM]}}
                \texttt{ПЕРЕСУНЬ: T, (1, 2)} \\
            }
        }
    }
    & & \\ \\

    \theutterance \stepcounter{utterance}  
    & & & \multicolumn{2}{p{0.3\linewidth}}{
        \cellcolor[rgb]{0.9,0.9,0.9}{
            \makecell[{{p{\linewidth}}}]{
                \texttt{\tiny{[GM$|$GM]}}
                \texttt{T moved to (1, 2)} \\
            }
        }
    }
    & & \\ \\

    \theutterance \stepcounter{utterance}  
    & & \multicolumn{4}{p{0.6\linewidth}}{
        \cellcolor[rgb]{0.9,0.9,0.9}{
            \makecell[{{p{\linewidth}}}]{
                \texttt{\tiny{[P1$\langle$GM]}}
                \texttt{Об'єкт 'M' успішно переміщено в (1,3). Ваша оновлена сітка виглядає так:} \\
\texttt{```} \\
\texttt{ 1234567} \\
\texttt{╔═══════╗} \\
\texttt{║TAD◌◌◌◌║ 1} \\
\texttt{║OPS◌◌◌◌║ 2} \\
\texttt{║M◌◌◌◌┌─╢ 3} \\
\texttt{╟──┐◌◌│◌║ 4} \\
\texttt{║◌◌├──┼─╢ 5} \\
\texttt{║◌◌│◌◌│◌║ 6} \\
\texttt{║◌◌│◌◌│◌║ 7} \\
\texttt{╚══╧══╧═╝} \\
\texttt{```} \\
\texttt{Зараз ви граєте раунд 8 з максимум 28.} \\
\texttt{Разом ви накопичили 6 з 16 штрафів.} \\
\texttt{Інший гравець перемістив об'єкт на своїй сітці.} \\
\\ 
\\ 
\texttt{Яка ваша наступна команда?} \\
            }
        }
    }
    & \\ \\

    \theutterance \stepcounter{utterance}  
    & \multicolumn{4}{p{0.6\linewidth}}{
        \cellcolor[rgb]{0.9,0.9,0.9}{
            \makecell[{{p{\linewidth}}}]{
                \texttt{\tiny{[P1$\rangle$GM]}}
                \texttt{СКАЖИ: завершено?} \\
            }
        }
    }
    & & \\ \\

    \theutterance \stepcounter{utterance}  
    & & & \multicolumn{2}{p{0.3\linewidth}}{
        \cellcolor[rgb]{0.9,0.9,0.9}{
            \makecell[{{p{\linewidth}}}]{
                \texttt{\tiny{[GM$|$GM]}}
                \texttt{СКАЖИ: завершено?} \\
            }
        }
    }
    & & \\ \\

    \theutterance \stepcounter{utterance}  
    & & \multicolumn{4}{p{0.6\linewidth}}{
        \cellcolor[rgb]{0.9,0.9,0.9}{
            \makecell[{{p{\linewidth}}}]{
                \texttt{\tiny{[GM$\rangle$P2]}}
                \texttt{Об'єкт 'T' успішно переміщено в (1,2). Ваша оновлена сітка виглядає так:} \\
\texttt{```} \\
\texttt{ 1234567} \\
\texttt{╔═══════╗} \\
\texttt{║PAD◌◌◌◌║ 1} \\
\texttt{║T◌S◌◌◌◌║ 2} \\
\texttt{║OM◌◌◌┌─╢ 3} \\
\texttt{╟──┐◌◌│◌║ 4} \\
\texttt{║◌◌├──┼─╢ 5} \\
\texttt{║◌◌│◌◌│◌║ 6} \\
\texttt{║◌◌│◌◌│◌║ 7} \\
\texttt{╚══╧══╧═╝} \\
\texttt{```} \\
\texttt{Зараз ви граєте раунд 9 з максимум 28.} \\
\texttt{Разом ви накопичили 6 з 16 штрафів.} \\
\texttt{Інший гравець надіслав це повідомлення:} \\
\\ 
\\ 
\texttt{"завершено?"} \\
\\ 
\\ 
\texttt{Яка ваша наступна команда?} \\
            }
        }
    }
    & \\ \\

    \theutterance \stepcounter{utterance}  
    & & & \multicolumn{4}{p{0.6\linewidth}}{
        \cellcolor[rgb]{0.9,0.9,0.9}{
            \makecell[{{p{\linewidth}}}]{
                \texttt{\tiny{[GM$\langle$P2]}}
                \texttt{СКАЖИ: ще не готовий, продовжуємо переміщувати фігури} \\
            }
        }
    }
     \\ \\

    \theutterance \stepcounter{utterance}  
    & & & \multicolumn{2}{p{0.3\linewidth}}{
        \cellcolor[rgb]{0.9,0.9,0.9}{
            \makecell[{{p{\linewidth}}}]{
                \texttt{\tiny{[GM$|$GM]}}
                \texttt{СКАЖИ: ще не готовий, продовжуємо переміщувати фігури} \\
            }
        }
    }
    & & \\ \\

    \theutterance \stepcounter{utterance}  
    & & \multicolumn{4}{p{0.6\linewidth}}{
        \cellcolor[rgb]{0.9,0.9,0.9}{
            \makecell[{{p{\linewidth}}}]{
                \texttt{\tiny{[P1$\langle$GM]}}
                \texttt{Ваше повідомлення передано іншому гравцю.} \\
\\ 
\\ 
\texttt{Зараз ви граєте раунд 9 з максимум 28.} \\
\texttt{Разом ви накопичили 6 з 16 штрафів.} \\
\texttt{Інший гравець надіслав це повідомлення:} \\
\\ 
\\ 
\texttt{"ще не готовий, продовжуємо переміщувати фігури"} \\
\\ 
\\ 
\texttt{Яка ваша наступна команда?} \\
            }
        }
    }
    & \\ \\

    \theutterance \stepcounter{utterance}  
    & \multicolumn{4}{p{0.6\linewidth}}{
        \cellcolor[rgb]{0.9,0.9,0.9}{
            \makecell[{{p{\linewidth}}}]{
                \texttt{\tiny{[P1$\rangle$GM]}}
                \texttt{СКАЖИ: Зрозумів. Будь ласка, повідомте, коли будете готові, щоб ми могли завершити гру одночасно.} \\
            }
        }
    }
    & & \\ \\

    \theutterance \stepcounter{utterance}  
    & & & \multicolumn{2}{p{0.3\linewidth}}{
        \cellcolor[rgb]{0.9,0.9,0.9}{
            \makecell[{{p{\linewidth}}}]{
                \texttt{\tiny{[GM$|$GM]}}
                \texttt{СКАЖИ: Зрозумів. Будь ласка, повідомте, коли будете готові, щоб ми могли завершити гру одночасно.} \\
            }
        }
    }
    & & \\ \\

    \theutterance \stepcounter{utterance}  
    & & \multicolumn{4}{p{0.6\linewidth}}{
        \cellcolor[rgb]{0.9,0.9,0.9}{
            \makecell[{{p{\linewidth}}}]{
                \texttt{\tiny{[GM$\rangle$P2]}}
                \texttt{Ваше повідомлення передано іншому гравцю.} \\
\\ 
\\ 
\texttt{Зараз ви граєте раунд 10 з максимум 28.} \\
\texttt{Разом ви накопичили 6 з 16 штрафів.} \\
\texttt{Інший гравець надіслав це повідомлення:} \\
\\ 
\\ 
\texttt{"Зрозумів. Будь ласка, повідомте, коли будете готові, щоб ми могли завершити гру одночасно."} \\
\\ 
\\ 
\texttt{Яка ваша наступна команда?} \\
            }
        }
    }
    & \\ \\

    \theutterance \stepcounter{utterance}  
    & & & \multicolumn{4}{p{0.6\linewidth}}{
        \cellcolor[rgb]{0.9,0.9,0.9}{
            \makecell[{{p{\linewidth}}}]{
                \texttt{\tiny{[GM$\langle$P2]}}
                \texttt{ПЕРЕСУНЬ: T, (1, 1)} \\
            }
        }
    }
     \\ \\

    \theutterance \stepcounter{utterance}  
    & & & \multicolumn{2}{p{0.3\linewidth}}{
        \cellcolor[rgb]{0.9,0.9,0.9}{
            \makecell[{{p{\linewidth}}}]{
                \texttt{\tiny{[GM$|$GM]}}
                \texttt{ПЕРЕСУНЬ: T, (1, 1)} \\
            }
        }
    }
    & & \\ \\

    \theutterance \stepcounter{utterance}  
    & & & \multicolumn{2}{p{0.3\linewidth}}{
        \cellcolor[rgb]{0.9,0.9,0.9}{
            \makecell[{{p{\linewidth}}}]{
                \texttt{\tiny{[GM$|$GM]}}
                \texttt{Штраф: (1,1) не порожня, а містить 'P'.} \\
\\ 
\texttt{Разом ви накопичили 7 з 16 штрафів.} \\
\texttt{Будь ласка, спробуйте ще раз!} \\
            }
        }
    }
    & & \\ \\

    \theutterance \stepcounter{utterance}  
    & & \multicolumn{4}{p{0.6\linewidth}}{
        \cellcolor[rgb]{0.9,0.9,0.9}{
            \makecell[{{p{\linewidth}}}]{
                \texttt{\tiny{[GM$\rangle$P2]}}
                \texttt{Штраф: (1,1) не порожня, а містить 'P'.} \\
\\ 
\texttt{Разом ви накопичили 7 з 16 штрафів.} \\
\texttt{Будь ласка, спробуйте ще раз!} \\
            }
        }
    }
    & \\ \\

    \theutterance \stepcounter{utterance}  
    & & & \multicolumn{4}{p{0.6\linewidth}}{
        \cellcolor[rgb]{0.9,0.9,0.9}{
            \makecell[{{p{\linewidth}}}]{
                \texttt{\tiny{[GM$\langle$P2]}}
                \texttt{ПЕРЕСУНЬ: T, (2, 1)} \\
            }
        }
    }
     \\ \\

    \theutterance \stepcounter{utterance}  
    & & & \multicolumn{2}{p{0.3\linewidth}}{
        \cellcolor[rgb]{0.9,0.9,0.9}{
            \makecell[{{p{\linewidth}}}]{
                \texttt{\tiny{[GM$|$GM]}}
                \texttt{ПЕРЕСУНЬ: T, (2, 1)} \\
            }
        }
    }
    & & \\ \\

    \theutterance \stepcounter{utterance}  
    & & & \multicolumn{2}{p{0.3\linewidth}}{
        \cellcolor[rgb]{0.9,0.9,0.9}{
            \makecell[{{p{\linewidth}}}]{
                \texttt{\tiny{[GM$|$GM]}}
                \texttt{Штраф: (2,1) не порожня, а містить 'A'.} \\
\\ 
\texttt{Разом ви накопичили 8 з 16 штрафів.} \\
\texttt{Будь ласка, спробуйте ще раз!} \\
            }
        }
    }
    & & \\ \\

    \theutterance \stepcounter{utterance}  
    & & \multicolumn{4}{p{0.6\linewidth}}{
        \cellcolor[rgb]{0.9,0.9,0.9}{
            \makecell[{{p{\linewidth}}}]{
                \texttt{\tiny{[GM$\rangle$P2]}}
                \texttt{Штраф: (2,1) не порожня, а містить 'A'.} \\
\\ 
\texttt{Разом ви накопичили 8 з 16 штрафів.} \\
\texttt{Будь ласка, спробуйте ще раз!} \\
            }
        }
    }
    & \\ \\

    \theutterance \stepcounter{utterance}  
    & & & \multicolumn{4}{p{0.6\linewidth}}{
        \cellcolor[rgb]{0.9,0.9,0.9}{
            \makecell[{{p{\linewidth}}}]{
                \texttt{\tiny{[GM$\langle$P2]}}
                \texttt{ПЕРЕСУНЬ: T, (3, 1)} \\
            }
        }
    }
     \\ \\

    \theutterance \stepcounter{utterance}  
    & & & \multicolumn{2}{p{0.3\linewidth}}{
        \cellcolor[rgb]{0.9,0.9,0.9}{
            \makecell[{{p{\linewidth}}}]{
                \texttt{\tiny{[GM$|$GM]}}
                \texttt{ПЕРЕСУНЬ: T, (3, 1)} \\
            }
        }
    }
    & & \\ \\

    \theutterance \stepcounter{utterance}  
    & & & \multicolumn{2}{p{0.3\linewidth}}{
        \cellcolor[rgb]{0.9,0.9,0.9}{
            \makecell[{{p{\linewidth}}}]{
                \texttt{\tiny{[GM$|$GM]}}
                \texttt{Штраф: (3,1) не порожня, а містить 'D'.} \\
\\ 
\texttt{Разом ви накопичили 9 з 16 штрафів.} \\
\texttt{Будь ласка, спробуйте ще раз!} \\
            }
        }
    }
    & & \\ \\

    \theutterance \stepcounter{utterance}  
    & & \multicolumn{4}{p{0.6\linewidth}}{
        \cellcolor[rgb]{0.9,0.9,0.9}{
            \makecell[{{p{\linewidth}}}]{
                \texttt{\tiny{[GM$\rangle$P2]}}
                \texttt{Штраф: (3,1) не порожня, а містить 'D'.} \\
\\ 
\texttt{Разом ви накопичили 9 з 16 штрафів.} \\
\texttt{Будь ласка, спробуйте ще раз!} \\
            }
        }
    }
    & \\ \\

    \theutterance \stepcounter{utterance}  
    & & & \multicolumn{4}{p{0.6\linewidth}}{
        \cellcolor[rgb]{0.9,0.9,0.9}{
            \makecell[{{p{\linewidth}}}]{
                \texttt{\tiny{[GM$\langle$P2]}}
                \texttt{ПЕРЕСУНЬ: T, (1, 3)} \\
            }
        }
    }
     \\ \\

    \theutterance \stepcounter{utterance}  
    & & & \multicolumn{2}{p{0.3\linewidth}}{
        \cellcolor[rgb]{0.9,0.9,0.9}{
            \makecell[{{p{\linewidth}}}]{
                \texttt{\tiny{[GM$|$GM]}}
                \texttt{ПЕРЕСУНЬ: T, (1, 3)} \\
            }
        }
    }
    & & \\ \\

    \theutterance \stepcounter{utterance}  
    & & & \multicolumn{2}{p{0.3\linewidth}}{
        \cellcolor[rgb]{0.9,0.9,0.9}{
            \makecell[{{p{\linewidth}}}]{
                \texttt{\tiny{[GM$|$GM]}}
                \texttt{Штраф: (1,3) не порожня, а містить 'O'.} \\
\\ 
\texttt{Разом ви накопичили 10 з 16 штрафів.} \\
\texttt{Будь ласка, спробуйте ще раз!} \\
            }
        }
    }
    & & \\ \\

    \theutterance \stepcounter{utterance}  
    & & \multicolumn{4}{p{0.6\linewidth}}{
        \cellcolor[rgb]{0.9,0.9,0.9}{
            \makecell[{{p{\linewidth}}}]{
                \texttt{\tiny{[GM$\rangle$P2]}}
                \texttt{Штраф: (1,3) не порожня, а містить 'O'.} \\
\\ 
\texttt{Разом ви накопичили 10 з 16 штрафів.} \\
\texttt{Будь ласка, спробуйте ще раз!} \\
            }
        }
    }
    & \\ \\

    \theutterance \stepcounter{utterance}  
    & & & \multicolumn{4}{p{0.6\linewidth}}{
        \cellcolor[rgb]{0.9,0.9,0.9}{
            \makecell[{{p{\linewidth}}}]{
                \texttt{\tiny{[GM$\langle$P2]}}
                \texttt{ПЕРЕСУНЬ: T, (2, 2)} \\
            }
        }
    }
     \\ \\

    \theutterance \stepcounter{utterance}  
    & & & \multicolumn{2}{p{0.3\linewidth}}{
        \cellcolor[rgb]{0.9,0.9,0.9}{
            \makecell[{{p{\linewidth}}}]{
                \texttt{\tiny{[GM$|$GM]}}
                \texttt{ПЕРЕСУНЬ: T, (2, 2)} \\
            }
        }
    }
    & & \\ \\

    \theutterance \stepcounter{utterance}  
    & & & \multicolumn{2}{p{0.3\linewidth}}{
        \cellcolor[rgb]{0.9,0.9,0.9}{
            \makecell[{{p{\linewidth}}}]{
                \texttt{\tiny{[GM$|$GM]}}
                \texttt{T moved to (2, 2)} \\
            }
        }
    }
    & & \\ \\

    \theutterance \stepcounter{utterance}  
    & & \multicolumn{4}{p{0.6\linewidth}}{
        \cellcolor[rgb]{0.9,0.9,0.9}{
            \makecell[{{p{\linewidth}}}]{
                \texttt{\tiny{[P1$\langle$GM]}}
                \texttt{Ваше повідомлення передано іншому гравцю.} \\
\\ 
\\ 
\texttt{Зараз ви граєте раунд 10 з максимум 28.} \\
\texttt{Разом ви накопичили 10 з 16 штрафів.} \\
\texttt{Інший гравець перемістив об'єкт на своїй сітці.} \\
\\ 
\\ 
\texttt{Яка ваша наступна команда?} \\
            }
        }
    }
    & \\ \\

    \theutterance \stepcounter{utterance}  
    & \multicolumn{4}{p{0.6\linewidth}}{
        \cellcolor[rgb]{0.9,0.9,0.9}{
            \makecell[{{p{\linewidth}}}]{
                \texttt{\tiny{[P1$\rangle$GM]}}
                \texttt{СКАЖИ: Я вже перемістив усі фігури на свої місця. Будь ласка, завершуйте свої переміщення і надішліть "завершено!", коли будете готові.} \\
            }
        }
    }
    & & \\ \\

    \theutterance \stepcounter{utterance}  
    & & & \multicolumn{2}{p{0.3\linewidth}}{
        \cellcolor[rgb]{0.9,0.9,0.9}{
            \makecell[{{p{\linewidth}}}]{
                \texttt{\tiny{[GM$|$GM]}}
                \texttt{СКАЖИ: Я вже перемістив усі фігури на свої місця. Будь ласка, завершуйте свої переміщення і надішліть "завершено!", коли будете готові.} \\
            }
        }
    }
    & & \\ \\

    \theutterance \stepcounter{utterance}  
    & & \multicolumn{4}{p{0.6\linewidth}}{
        \cellcolor[rgb]{0.9,0.9,0.9}{
            \makecell[{{p{\linewidth}}}]{
                \texttt{\tiny{[GM$\rangle$P2]}}
                \texttt{Об'єкт 'T' успішно переміщено в (2,2). Ваша оновлена сітка виглядає так:} \\
\texttt{```} \\
\texttt{ 1234567} \\
\texttt{╔═══════╗} \\
\texttt{║PAD◌◌◌◌║ 1} \\
\texttt{║◌TS◌◌◌◌║ 2} \\
\texttt{║OM◌◌◌┌─╢ 3} \\
\texttt{╟──┐◌◌│◌║ 4} \\
\texttt{║◌◌├──┼─╢ 5} \\
\texttt{║◌◌│◌◌│◌║ 6} \\
\texttt{║◌◌│◌◌│◌║ 7} \\
\texttt{╚══╧══╧═╝} \\
\texttt{```} \\
\texttt{Зараз ви граєте раунд 11 з максимум 28.} \\
\texttt{Разом ви накопичили 10 з 16 штрафів.} \\
\texttt{Інший гравець надіслав це повідомлення:} \\
\\ 
\\ 
\texttt{"Я вже перемістив усі фігури на свої місця. Будь ласка, завершуйте свої переміщення і надішліть "завершено!", коли будете готові."} \\
\\ 
\\ 
\texttt{Яка ваша наступна команда?} \\
            }
        }
    }
    & \\ \\

    \theutterance \stepcounter{utterance}  
    & & & \multicolumn{4}{p{0.6\linewidth}}{
        \cellcolor[rgb]{0.9,0.9,0.9}{
            \makecell[{{p{\linewidth}}}]{
                \texttt{\tiny{[GM$\langle$P2]}}
                \texttt{ПЕРЕСУНЬ: T, (1, 1)} \\
            }
        }
    }
     \\ \\

    \theutterance \stepcounter{utterance}  
    & & & \multicolumn{2}{p{0.3\linewidth}}{
        \cellcolor[rgb]{0.9,0.9,0.9}{
            \makecell[{{p{\linewidth}}}]{
                \texttt{\tiny{[GM$|$GM]}}
                \texttt{ПЕРЕСУНЬ: T, (1, 1)} \\
            }
        }
    }
    & & \\ \\

    \theutterance \stepcounter{utterance}  
    & & & \multicolumn{2}{p{0.3\linewidth}}{
        \cellcolor[rgb]{0.9,0.9,0.9}{
            \makecell[{{p{\linewidth}}}]{
                \texttt{\tiny{[GM$|$GM]}}
                \texttt{Штраф: (1,1) не порожня, а містить 'P'.} \\
\\ 
\texttt{Разом ви накопичили 11 з 16 штрафів.} \\
\texttt{Будь ласка, спробуйте ще раз!} \\
            }
        }
    }
    & & \\ \\

    \theutterance \stepcounter{utterance}  
    & & \multicolumn{4}{p{0.6\linewidth}}{
        \cellcolor[rgb]{0.9,0.9,0.9}{
            \makecell[{{p{\linewidth}}}]{
                \texttt{\tiny{[GM$\rangle$P2]}}
                \texttt{Штраф: (1,1) не порожня, а містить 'P'.} \\
\\ 
\texttt{Разом ви накопичили 11 з 16 штрафів.} \\
\texttt{Будь ласка, спробуйте ще раз!} \\
            }
        }
    }
    & \\ \\

    \theutterance \stepcounter{utterance}  
    & & & \multicolumn{4}{p{0.6\linewidth}}{
        \cellcolor[rgb]{0.9,0.9,0.9}{
            \makecell[{{p{\linewidth}}}]{
                \texttt{\tiny{[GM$\langle$P2]}}
                \texttt{ПЕРЕСУНЬ: T, (2, 1)} \\
            }
        }
    }
     \\ \\

    \theutterance \stepcounter{utterance}  
    & & & \multicolumn{2}{p{0.3\linewidth}}{
        \cellcolor[rgb]{0.9,0.9,0.9}{
            \makecell[{{p{\linewidth}}}]{
                \texttt{\tiny{[GM$|$GM]}}
                \texttt{ПЕРЕСУНЬ: T, (2, 1)} \\
            }
        }
    }
    & & \\ \\

    \theutterance \stepcounter{utterance}  
    & & & \multicolumn{2}{p{0.3\linewidth}}{
        \cellcolor[rgb]{0.9,0.9,0.9}{
            \makecell[{{p{\linewidth}}}]{
                \texttt{\tiny{[GM$|$GM]}}
                \texttt{Штраф: (2,1) не порожня, а містить 'A'.} \\
\\ 
\texttt{Разом ви накопичили 12 з 16 штрафів.} \\
\texttt{Будь ласка, спробуйте ще раз!} \\
            }
        }
    }
    & & \\ \\

    \theutterance \stepcounter{utterance}  
    & & \multicolumn{4}{p{0.6\linewidth}}{
        \cellcolor[rgb]{0.9,0.9,0.9}{
            \makecell[{{p{\linewidth}}}]{
                \texttt{\tiny{[GM$\rangle$P2]}}
                \texttt{Штраф: (2,1) не порожня, а містить 'A'.} \\
\\ 
\texttt{Разом ви накопичили 12 з 16 штрафів.} \\
\texttt{Будь ласка, спробуйте ще раз!} \\
            }
        }
    }
    & \\ \\

    \theutterance \stepcounter{utterance}  
    & & & \multicolumn{4}{p{0.6\linewidth}}{
        \cellcolor[rgb]{0.9,0.9,0.9}{
            \makecell[{{p{\linewidth}}}]{
                \texttt{\tiny{[GM$\langle$P2]}}
                \texttt{ПЕРЕСУНЬ: T, (3, 1)} \\
            }
        }
    }
     \\ \\

    \theutterance \stepcounter{utterance}  
    & & & \multicolumn{2}{p{0.3\linewidth}}{
        \cellcolor[rgb]{0.9,0.9,0.9}{
            \makecell[{{p{\linewidth}}}]{
                \texttt{\tiny{[GM$|$GM]}}
                \texttt{ПЕРЕСУНЬ: T, (3, 1)} \\
            }
        }
    }
    & & \\ \\

    \theutterance \stepcounter{utterance}  
    & & & \multicolumn{2}{p{0.3\linewidth}}{
        \cellcolor[rgb]{0.9,0.9,0.9}{
            \makecell[{{p{\linewidth}}}]{
                \texttt{\tiny{[GM$|$GM]}}
                \texttt{Штраф: (3,1) не порожня, а містить 'D'.} \\
\\ 
\texttt{Разом ви накопичили 13 з 16 штрафів.} \\
\texttt{Будь ласка, спробуйте ще раз!} \\
            }
        }
    }
    & & \\ \\

    \theutterance \stepcounter{utterance}  
    & & \multicolumn{4}{p{0.6\linewidth}}{
        \cellcolor[rgb]{0.9,0.9,0.9}{
            \makecell[{{p{\linewidth}}}]{
                \texttt{\tiny{[GM$\rangle$P2]}}
                \texttt{Штраф: (3,1) не порожня, а містить 'D'.} \\
\\ 
\texttt{Разом ви накопичили 13 з 16 штрафів.} \\
\texttt{Будь ласка, спробуйте ще раз!} \\
            }
        }
    }
    & \\ \\

    \theutterance \stepcounter{utterance}  
    & & & \multicolumn{4}{p{0.6\linewidth}}{
        \cellcolor[rgb]{0.9,0.9,0.9}{
            \makecell[{{p{\linewidth}}}]{
                \texttt{\tiny{[GM$\langle$P2]}}
                \texttt{ПЕРЕСУНЬ: T, (1, 2)} \\
            }
        }
    }
     \\ \\

    \theutterance \stepcounter{utterance}  
    & & & \multicolumn{2}{p{0.3\linewidth}}{
        \cellcolor[rgb]{0.9,0.9,0.9}{
            \makecell[{{p{\linewidth}}}]{
                \texttt{\tiny{[GM$|$GM]}}
                \texttt{ПЕРЕСУНЬ: T, (1, 2)} \\
            }
        }
    }
    & & \\ \\

    \theutterance \stepcounter{utterance}  
    & & & \multicolumn{2}{p{0.3\linewidth}}{
        \cellcolor[rgb]{0.9,0.9,0.9}{
            \makecell[{{p{\linewidth}}}]{
                \texttt{\tiny{[GM$|$GM]}}
                \texttt{T moved to (1, 2)} \\
            }
        }
    }
    & & \\ \\

    \theutterance \stepcounter{utterance}  
    & & \multicolumn{4}{p{0.6\linewidth}}{
        \cellcolor[rgb]{0.9,0.9,0.9}{
            \makecell[{{p{\linewidth}}}]{
                \texttt{\tiny{[P1$\langle$GM]}}
                \texttt{Ваше повідомлення передано іншому гравцю.} \\
\\ 
\\ 
\texttt{Зараз ви граєте раунд 11 з максимум 28.} \\
\texttt{Разом ви накопичили 13 з 16 штрафів.} \\
\texttt{Інший гравець перемістив об'єкт на своїй сітці.} \\
\\ 
\\ 
\texttt{Яка ваша наступна команда?} \\
            }
        }
    }
    & \\ \\

    \theutterance \stepcounter{utterance}  
    & \multicolumn{4}{p{0.6\linewidth}}{
        \cellcolor[rgb]{0.9,0.9,0.9}{
            \makecell[{{p{\linewidth}}}]{
                \texttt{\tiny{[P1$\rangle$GM]}}
                \texttt{СКАЖИ: зачекай, я перевірю свої координати. Т: (1,1), A: (2,1), D: (3,1), O: (1,2), P: (2,2), S: (3,2), M: (1,3). Все правильно?} \\
            }
        }
    }
    & & \\ \\

    \theutterance \stepcounter{utterance}  
    & & & \multicolumn{2}{p{0.3\linewidth}}{
        \cellcolor[rgb]{0.9,0.9,0.9}{
            \makecell[{{p{\linewidth}}}]{
                \texttt{\tiny{[GM$|$GM]}}
                \texttt{СКАЖИ: зачекай, я перевірю свої координати. Т: (1,1), A: (2,1), D: (3,1), O: (1,2), P: (2,2), S: (3,2), M: (1,3). Все правильно?} \\
            }
        }
    }
    & & \\ \\

    \theutterance \stepcounter{utterance}  
    & & \multicolumn{4}{p{0.6\linewidth}}{
        \cellcolor[rgb]{0.9,0.9,0.9}{
            \makecell[{{p{\linewidth}}}]{
                \texttt{\tiny{[GM$\rangle$P2]}}
                \texttt{Об'єкт 'T' успішно переміщено в (1,2). Ваша оновлена сітка виглядає так:} \\
\texttt{```} \\
\texttt{ 1234567} \\
\texttt{╔═══════╗} \\
\texttt{║PAD◌◌◌◌║ 1} \\
\texttt{║T◌S◌◌◌◌║ 2} \\
\texttt{║OM◌◌◌┌─╢ 3} \\
\texttt{╟──┐◌◌│◌║ 4} \\
\texttt{║◌◌├──┼─╢ 5} \\
\texttt{║◌◌│◌◌│◌║ 6} \\
\texttt{║◌◌│◌◌│◌║ 7} \\
\texttt{╚══╧══╧═╝} \\
\texttt{```} \\
\texttt{Зараз ви граєте раунд 12 з максимум 28.} \\
\texttt{Разом ви накопичили 13 з 16 штрафів.} \\
\texttt{Інший гравець надіслав це повідомлення:} \\
\\ 
\\ 
\texttt{"зачекай, я перевірю свої координати. Т: (1,1), A: (2,1), D: (3,1), O: (1,2), P: (2,2), S: (3,2), M: (1,3). Все правильно?"} \\
\\ 
\\ 
\texttt{Яка ваша наступна команда?} \\
            }
        }
    }
    & \\ \\

    \theutterance \stepcounter{utterance}  
    & & & \multicolumn{4}{p{0.6\linewidth}}{
        \cellcolor[rgb]{0.9,0.9,0.9}{
            \makecell[{{p{\linewidth}}}]{
                \texttt{\tiny{[GM$\langle$P2]}}
                \texttt{СКАЖИ: Ні, T має бути на (1,1). Зараз T на (1,2). Будь ласка, перемісти T на (1,1). P має бути на (2,2), S на (3,2), M на (1,3). Перевір, чи всі фігури на своїх місцях: P(2,2), A(2,1), D(3,1), O(1,2), S(3,2), M(1,3), T(1,1). У мене P(1,1), A(2,1), D(3,1), O(1,3), S(3,2), M(2,3), T(1,2). Тобто нам потрібно узгодити позиції. Моя цільова схема така сама, як твоя початкова. Тобто: T(1,1), A(2,1), D(3,1), O(1,2), P(2,2), S(3,2), M(1,3). Перевір, чи ти дійсно перемістив P на (2,2), S на (3,2), M на (1,3) і T на (1,1). Якщо так, то я переміщу свої фігури туди. Зараз я переміщую P на (2,2).} \\
            }
        }
    }
     \\ \\

    \theutterance \stepcounter{utterance}  
    & & & \multicolumn{2}{p{0.3\linewidth}}{
        \cellcolor[rgb]{0.9,0.9,0.9}{
            \makecell[{{p{\linewidth}}}]{
                \texttt{\tiny{[GM$|$GM]}}
                \texttt{СКАЖИ: Ні, T має бути на (1,1). Зараз T на (1,2). Будь ласка, перемісти T на (1,1). P має бути на (2,2), S на (3,2), M на (1,3). Перевір, чи всі фігури на своїх місцях: P(2,2), A(2,1), D(3,1), O(1,2), S(3,2), M(1,3), T(1,1). У мене P(1,1), A(2,1), D(3,1), O(1,3), S(3,2), M(2,3), T(1,2). Тобто нам потрібно узгодити позиції. Моя цільова схема така сама, як твоя початкова. Тобто: T(1,1), A(2,1), D(3,1), O(1,2), P(2,2), S(3,2), M(1,3). Перевір, чи ти дійсно перемістив P на (2,2), S на (3,2), M на (1,3) і T на (1,1). Якщо так, то я переміщу свої фігури туди. Зараз я переміщую P на (2,2).} \\
            }
        }
    }
    & & \\ \\

    \theutterance \stepcounter{utterance}  
    & & \multicolumn{4}{p{0.6\linewidth}}{
        \cellcolor[rgb]{0.9,0.9,0.9}{
            \makecell[{{p{\linewidth}}}]{
                \texttt{\tiny{[P1$\langle$GM]}}
                \texttt{Ваше повідомлення передано іншому гравцю.} \\
\\ 
\\ 
\texttt{Зараз ви граєте раунд 12 з максимум 28.} \\
\texttt{Разом ви накопичили 13 з 16 штрафів.} \\
\texttt{Інший гравець надіслав це повідомлення:} \\
\\ 
\\ 
\texttt{"Ні, T має бути на (1,1). Зараз T на (1,2). Будь ласка, перемісти T на (1,1). P має бути на (2,2), S на (3,2), M на (1,3). Перевір, чи всі фігури на своїх місцях: P(2,2), A(2,1), D(3,1), O(1,2), S(3,2), M(1,3), T(1,1). У мене P(1,1), A(2,1), D(3,1), O(1,3), S(3,2), M(2,3), T(1,2). Тобто нам потрібно узгодити позиції. Моя цільова схема така сама, як твоя початкова. Тобто: T(1,1), A(2,1), D(3,1), O(1,2), P(2,2), S(3,2), M(1,3). Перевір, чи ти дійсно перемістив P на (2,2), S на (3,2), M на (1,3) і T на (1,1). Якщо так, то я переміщу свої фігури туди. Зараз я переміщую P на (2,2)."} \\
\\ 
\\ 
\texttt{Яка ваша наступна команда?} \\
            }
        }
    }
    & \\ \\

    \theutterance \stepcounter{utterance}  
    & \multicolumn{4}{p{0.6\linewidth}}{
        \cellcolor[rgb]{0.9,0.9,0.9}{
            \makecell[{{p{\linewidth}}}]{
                \texttt{\tiny{[P1$\rangle$GM]}}
                \texttt{СКАЖИ: Так, я вже перемістив T, A, D, O, P, S, M на зазначені координати. Чекаю на твоє "завершено!".} \\
            }
        }
    }
    & & \\ \\

    \theutterance \stepcounter{utterance}  
    & & & \multicolumn{2}{p{0.3\linewidth}}{
        \cellcolor[rgb]{0.9,0.9,0.9}{
            \makecell[{{p{\linewidth}}}]{
                \texttt{\tiny{[GM$|$GM]}}
                \texttt{СКАЖИ: Так, я вже перемістив T, A, D, O, P, S, M на зазначені координати. Чекаю на твоє "завершено!".} \\
            }
        }
    }
    & & \\ \\

    \theutterance \stepcounter{utterance}  
    & & \multicolumn{4}{p{0.6\linewidth}}{
        \cellcolor[rgb]{0.9,0.9,0.9}{
            \makecell[{{p{\linewidth}}}]{
                \texttt{\tiny{[GM$\rangle$P2]}}
                \texttt{Ваше повідомлення передано іншому гравцю.} \\
\\ 
\\ 
\texttt{Зараз ви граєте раунд 13 з максимум 28.} \\
\texttt{Разом ви накопичили 13 з 16 штрафів.} \\
\texttt{Інший гравець надіслав це повідомлення:} \\
\\ 
\\ 
\texttt{"Так, я вже перемістив T, A, D, O, P, S, M на зазначені координати. Чекаю на твоє "завершено!"."} \\
\\ 
\\ 
\texttt{Яка ваша наступна команда?} \\
            }
        }
    }
    & \\ \\

    \theutterance \stepcounter{utterance}  
    & & & \multicolumn{4}{p{0.6\linewidth}}{
        \cellcolor[rgb]{0.9,0.9,0.9}{
            \makecell[{{p{\linewidth}}}]{
                \texttt{\tiny{[GM$\langle$P2]}}
                \texttt{СКАЖИ: зачекай, я ще не закінчив. Мені потрібно перемістити P на (2,2), S на (3,2), M на (1,3) і T на (1,1). Зараз я переміщую P.} \\
            }
        }
    }
     \\ \\

    \theutterance \stepcounter{utterance}  
    & & & \multicolumn{2}{p{0.3\linewidth}}{
        \cellcolor[rgb]{0.9,0.9,0.9}{
            \makecell[{{p{\linewidth}}}]{
                \texttt{\tiny{[GM$|$GM]}}
                \texttt{СКАЖИ: зачекай, я ще не закінчив. Мені потрібно перемістити P на (2,2), S на (3,2), M на (1,3) і T на (1,1). Зараз я переміщую P.} \\
            }
        }
    }
    & & \\ \\

    \theutterance \stepcounter{utterance}  
    & & \multicolumn{4}{p{0.6\linewidth}}{
        \cellcolor[rgb]{0.9,0.9,0.9}{
            \makecell[{{p{\linewidth}}}]{
                \texttt{\tiny{[P1$\langle$GM]}}
                \texttt{Ваше повідомлення передано іншому гравцю.} \\
\\ 
\\ 
\texttt{Зараз ви граєте раунд 13 з максимум 28.} \\
\texttt{Разом ви накопичили 13 з 16 штрафів.} \\
\texttt{Інший гравець надіслав це повідомлення:} \\
\\ 
\\ 
\texttt{"зачекай, я ще не закінчив. Мені потрібно перемістити P на (2,2), S на (3,2), M на (1,3) і T на (1,1). Зараз я переміщую P."} \\
\\ 
\\ 
\texttt{Яка ваша наступна команда?} \\
            }
        }
    }
    & \\ \\

    \theutterance \stepcounter{utterance}  
    & \multicolumn{4}{p{0.6\linewidth}}{
        \cellcolor[rgb]{0.9,0.9,0.9}{
            \makecell[{{p{\linewidth}}}]{
                \texttt{\tiny{[P1$\rangle$GM]}}
                \texttt{СКАЖИ: Добре, я чекаю. Напиши "завершено!", коли ти теж перемістиш усі фігури на свої місця.} \\
            }
        }
    }
    & & \\ \\

    \theutterance \stepcounter{utterance}  
    & & & \multicolumn{2}{p{0.3\linewidth}}{
        \cellcolor[rgb]{0.9,0.9,0.9}{
            \makecell[{{p{\linewidth}}}]{
                \texttt{\tiny{[GM$|$GM]}}
                \texttt{СКАЖИ: Добре, я чекаю. Напиши "завершено!", коли ти теж перемістиш усі фігури на свої місця.} \\
            }
        }
    }
    & & \\ \\

    \theutterance \stepcounter{utterance}  
    & & \multicolumn{4}{p{0.6\linewidth}}{
        \cellcolor[rgb]{0.9,0.9,0.9}{
            \makecell[{{p{\linewidth}}}]{
                \texttt{\tiny{[GM$\rangle$P2]}}
                \texttt{Ваше повідомлення передано іншому гравцю.} \\
\\ 
\\ 
\texttt{Зараз ви граєте раунд 14 з максимум 28.} \\
\texttt{Разом ви накопичили 13 з 16 штрафів.} \\
\texttt{Інший гравець надіслав це повідомлення:} \\
\\ 
\\ 
\texttt{"Добре, я чекаю. Напиши "завершено!", коли ти теж перемістиш усі фігури на свої місця."} \\
\\ 
\\ 
\texttt{Яка ваша наступна команда?} \\
            }
        }
    }
    & \\ \\

    \theutterance \stepcounter{utterance}  
    & & & \multicolumn{4}{p{0.6\linewidth}}{
        \cellcolor[rgb]{0.9,0.9,0.9}{
            \makecell[{{p{\linewidth}}}]{
                \texttt{\tiny{[GM$\langle$P2]}}
                \texttt{ПЕРЕСУНЬ: P, (2, 2)} \\
            }
        }
    }
     \\ \\

    \theutterance \stepcounter{utterance}  
    & & & \multicolumn{2}{p{0.3\linewidth}}{
        \cellcolor[rgb]{0.9,0.9,0.9}{
            \makecell[{{p{\linewidth}}}]{
                \texttt{\tiny{[GM$|$GM]}}
                \texttt{ПЕРЕСУНЬ: P, (2, 2)} \\
            }
        }
    }
    & & \\ \\

    \theutterance \stepcounter{utterance}  
    & & & \multicolumn{2}{p{0.3\linewidth}}{
        \cellcolor[rgb]{0.9,0.9,0.9}{
            \makecell[{{p{\linewidth}}}]{
                \texttt{\tiny{[GM$|$GM]}}
                \texttt{P moved to (2, 2)} \\
            }
        }
    }
    & & \\ \\

    \theutterance \stepcounter{utterance}  
    & & \multicolumn{4}{p{0.6\linewidth}}{
        \cellcolor[rgb]{0.9,0.9,0.9}{
            \makecell[{{p{\linewidth}}}]{
                \texttt{\tiny{[P1$\langle$GM]}}
                \texttt{Ваше повідомлення передано іншому гравцю.} \\
\\ 
\\ 
\texttt{Зараз ви граєте раунд 14 з максимум 28.} \\
\texttt{Разом ви накопичили 13 з 16 штрафів.} \\
\texttt{Інший гравець перемістив об'єкт на своїй сітці.} \\
\\ 
\\ 
\texttt{Яка ваша наступна команда?} \\
            }
        }
    }
    & \\ \\

    \theutterance \stepcounter{utterance}  
    & \multicolumn{4}{p{0.6\linewidth}}{
        \cellcolor[rgb]{0.9,0.9,0.9}{
            \makecell[{{p{\linewidth}}}]{
                \texttt{\tiny{[P1$\rangle$GM]}}
                \texttt{СКАЖИ: зачекай, я перевірю свої координати. Т: (1,1), A: (2,1), D: (3,1), O: (1,2), P: (2,2), S: (3,2), M: (1,3). Все правильно?} \\
            }
        }
    }
    & & \\ \\

    \theutterance \stepcounter{utterance}  
    & & & \multicolumn{2}{p{0.3\linewidth}}{
        \cellcolor[rgb]{0.9,0.9,0.9}{
            \makecell[{{p{\linewidth}}}]{
                \texttt{\tiny{[GM$|$GM]}}
                \texttt{СКАЖИ: зачекай, я перевірю свої координати. Т: (1,1), A: (2,1), D: (3,1), O: (1,2), P: (2,2), S: (3,2), M: (1,3). Все правильно?} \\
            }
        }
    }
    & & \\ \\

    \theutterance \stepcounter{utterance}  
    & & \multicolumn{4}{p{0.6\linewidth}}{
        \cellcolor[rgb]{0.9,0.9,0.9}{
            \makecell[{{p{\linewidth}}}]{
                \texttt{\tiny{[GM$\rangle$P2]}}
                \texttt{Об'єкт 'P' успішно переміщено в (2,2). Ваша оновлена сітка виглядає так:} \\
\texttt{```} \\
\texttt{ 1234567} \\
\texttt{╔═══════╗} \\
\texttt{║◌AD◌◌◌◌║ 1} \\
\texttt{║TPS◌◌◌◌║ 2} \\
\texttt{║OM◌◌◌┌─╢ 3} \\
\texttt{╟──┐◌◌│◌║ 4} \\
\texttt{║◌◌├──┼─╢ 5} \\
\texttt{║◌◌│◌◌│◌║ 6} \\
\texttt{║◌◌│◌◌│◌║ 7} \\
\texttt{╚══╧══╧═╝} \\
\texttt{```} \\
\texttt{Зараз ви граєте раунд 15 з максимум 28.} \\
\texttt{Разом ви накопичили 13 з 16 штрафів.} \\
\texttt{Інший гравець надіслав це повідомлення:} \\
\\ 
\\ 
\texttt{"зачекай, я перевірю свої координати. Т: (1,1), A: (2,1), D: (3,1), O: (1,2), P: (2,2), S: (3,2), M: (1,3). Все правильно?"} \\
\\ 
\\ 
\texttt{Яка ваша наступна команда?} \\
            }
        }
    }
    & \\ \\

    \theutterance \stepcounter{utterance}  
    & & & \multicolumn{4}{p{0.6\linewidth}}{
        \cellcolor[rgb]{0.9,0.9,0.9}{
            \makecell[{{p{\linewidth}}}]{
                \texttt{\tiny{[GM$\langle$P2]}}
                \texttt{СКАЖИ: Ні, твоя цільова схема була: T(1,1), A(2,1), D(3,1), O(1,2), P(2,2), S(3,2), M(1,3). Але на твоєму полі зараз T(1,1), A(2,1), D(3,1), O(1,2), P(2,2), S(3,2), M(1,3). Це збігається з твоєю заявою. Але у мене на полі: P(2,2), A(2,1), D(3,1), T(1,2), S(3,2), O(1,3), M(2,3). Мені потрібно перемістити T на (1,1), O на (1,2), M на (1,3). Зараз я переміщую T на (1,1).} \\
            }
        }
    }
     \\ \\

    \theutterance \stepcounter{utterance}  
    & & & \multicolumn{2}{p{0.3\linewidth}}{
        \cellcolor[rgb]{0.9,0.9,0.9}{
            \makecell[{{p{\linewidth}}}]{
                \texttt{\tiny{[GM$|$GM]}}
                \texttt{СКАЖИ: Ні, твоя цільова схема була: T(1,1), A(2,1), D(3,1), O(1,2), P(2,2), S(3,2), M(1,3). Але на твоєму полі зараз T(1,1), A(2,1), D(3,1), O(1,2), P(2,2), S(3,2), M(1,3). Це збігається з твоєю заявою. Але у мене на полі: P(2,2), A(2,1), D(3,1), T(1,2), S(3,2), O(1,3), M(2,3). Мені потрібно перемістити T на (1,1), O на (1,2), M на (1,3). Зараз я переміщую T на (1,1).} \\
            }
        }
    }
    & & \\ \\

    \theutterance \stepcounter{utterance}  
    & & \multicolumn{4}{p{0.6\linewidth}}{
        \cellcolor[rgb]{0.9,0.9,0.9}{
            \makecell[{{p{\linewidth}}}]{
                \texttt{\tiny{[P1$\langle$GM]}}
                \texttt{Ваше повідомлення передано іншому гравцю.} \\
\\ 
\\ 
\texttt{Зараз ви граєте раунд 15 з максимум 28.} \\
\texttt{Разом ви накопичили 13 з 16 штрафів.} \\
\texttt{Інший гравець надіслав це повідомлення:} \\
\\ 
\\ 
\texttt{"Ні, твоя цільова схема була: T(1,1), A(2,1), D(3,1), O(1,2), P(2,2), S(3,2), M(1,3). Але на твоєму полі зараз T(1,1), A(2,1), D(3,1), O(1,2), P(2,2), S(3,2), M(1,3). Це збігається з твоєю заявою. Але у мене на полі: P(2,2), A(2,1), D(3,1), T(1,2), S(3,2), O(1,3), M(2,3). Мені потрібно перемістити T на (1,1), O на (1,2), M на (1,3). Зараз я переміщую T на (1,1)."} \\
\\ 
\\ 
\texttt{Яка ваша наступна команда?} \\
            }
        }
    }
    & \\ \\

    \theutterance \stepcounter{utterance}  
    & \multicolumn{4}{p{0.6\linewidth}}{
        \cellcolor[rgb]{0.9,0.9,0.9}{
            \makecell[{{p{\linewidth}}}]{
                \texttt{\tiny{[P1$\rangle$GM]}}
                \texttt{СКАЖИ: Чудово, ми на фінішній прямій. Я вже на місці. Напиши "завершено!", коли ти перемістиш T, O та M на свої місця.} \\
            }
        }
    }
    & & \\ \\

    \theutterance \stepcounter{utterance}  
    & & & \multicolumn{2}{p{0.3\linewidth}}{
        \cellcolor[rgb]{0.9,0.9,0.9}{
            \makecell[{{p{\linewidth}}}]{
                \texttt{\tiny{[GM$|$GM]}}
                \texttt{СКАЖИ: Чудово, ми на фінішній прямій. Я вже на місці. Напиши "завершено!", коли ти перемістиш T, O та M на свої місця.} \\
            }
        }
    }
    & & \\ \\

    \theutterance \stepcounter{utterance}  
    & & \multicolumn{4}{p{0.6\linewidth}}{
        \cellcolor[rgb]{0.9,0.9,0.9}{
            \makecell[{{p{\linewidth}}}]{
                \texttt{\tiny{[GM$\rangle$P2]}}
                \texttt{Ваше повідомлення передано іншому гравцю.} \\
\\ 
\\ 
\texttt{Зараз ви граєте раунд 16 з максимум 28.} \\
\texttt{Разом ви накопичили 13 з 16 штрафів.} \\
\texttt{Інший гравець надіслав це повідомлення:} \\
\\ 
\\ 
\texttt{"Чудово, ми на фінішній прямій. Я вже на місці. Напиши "завершено!", коли ти перемістиш T, O та M на свої місця."} \\
\\ 
\\ 
\texttt{Яка ваша наступна команда?} \\
            }
        }
    }
    & \\ \\

    \theutterance \stepcounter{utterance}  
    & & & \multicolumn{4}{p{0.6\linewidth}}{
        \cellcolor[rgb]{0.9,0.9,0.9}{
            \makecell[{{p{\linewidth}}}]{
                \texttt{\tiny{[GM$\langle$P2]}}
                \texttt{ПЕРЕСУНЬ: T, (1, 1)} \\
            }
        }
    }
     \\ \\

    \theutterance \stepcounter{utterance}  
    & & & \multicolumn{2}{p{0.3\linewidth}}{
        \cellcolor[rgb]{0.9,0.9,0.9}{
            \makecell[{{p{\linewidth}}}]{
                \texttt{\tiny{[GM$|$GM]}}
                \texttt{ПЕРЕСУНЬ: T, (1, 1)} \\
            }
        }
    }
    & & \\ \\

    \theutterance \stepcounter{utterance}  
    & & & \multicolumn{2}{p{0.3\linewidth}}{
        \cellcolor[rgb]{0.9,0.9,0.9}{
            \makecell[{{p{\linewidth}}}]{
                \texttt{\tiny{[GM$|$GM]}}
                \texttt{T moved to (1, 1)} \\
            }
        }
    }
    & & \\ \\

    \theutterance \stepcounter{utterance}  
    & & \multicolumn{4}{p{0.6\linewidth}}{
        \cellcolor[rgb]{0.9,0.9,0.9}{
            \makecell[{{p{\linewidth}}}]{
                \texttt{\tiny{[P1$\langle$GM]}}
                \texttt{Ваше повідомлення передано іншому гравцю.} \\
\\ 
\\ 
\texttt{Зараз ви граєте раунд 16 з максимум 28.} \\
\texttt{Разом ви накопичили 13 з 16 штрафів.} \\
\texttt{Інший гравець перемістив об'єкт на своїй сітці.} \\
\\ 
\\ 
\texttt{Яка ваша наступна команда?} \\
            }
        }
    }
    & \\ \\

    \theutterance \stepcounter{utterance}  
    & \multicolumn{4}{p{0.6\linewidth}}{
        \cellcolor[rgb]{0.9,0.9,0.9}{
            \makecell[{{p{\linewidth}}}]{
                \texttt{\tiny{[P1$\rangle$GM]}}
                \texttt{СКАЖИ: зачекай, я перевірю свої координати. Т: (1,1), A: (2,1), D: (3,1), O: (1,2), P: (2,2), S: (3,2), M: (1,3). Все правильно?} \\
            }
        }
    }
    & & \\ \\

    \theutterance \stepcounter{utterance}  
    & & & \multicolumn{2}{p{0.3\linewidth}}{
        \cellcolor[rgb]{0.9,0.9,0.9}{
            \makecell[{{p{\linewidth}}}]{
                \texttt{\tiny{[GM$|$GM]}}
                \texttt{СКАЖИ: зачекай, я перевірю свої координати. Т: (1,1), A: (2,1), D: (3,1), O: (1,2), P: (2,2), S: (3,2), M: (1,3). Все правильно?} \\
            }
        }
    }
    & & \\ \\

    \theutterance \stepcounter{utterance}  
    & & \multicolumn{4}{p{0.6\linewidth}}{
        \cellcolor[rgb]{0.9,0.9,0.9}{
            \makecell[{{p{\linewidth}}}]{
                \texttt{\tiny{[GM$\rangle$P2]}}
                \texttt{Об'єкт 'T' успішно переміщено в (1,1). Ваша оновлена сітка виглядає так:} \\
\texttt{```} \\
\texttt{ 1234567} \\
\texttt{╔═══════╗} \\
\texttt{║TAD◌◌◌◌║ 1} \\
\texttt{║◌PS◌◌◌◌║ 2} \\
\texttt{║OM◌◌◌┌─╢ 3} \\
\texttt{╟──┐◌◌│◌║ 4} \\
\texttt{║◌◌├──┼─╢ 5} \\
\texttt{║◌◌│◌◌│◌║ 6} \\
\texttt{║◌◌│◌◌│◌║ 7} \\
\texttt{╚══╧══╧═╝} \\
\texttt{```} \\
\texttt{Зараз ви граєте раунд 17 з максимум 28.} \\
\texttt{Разом ви накопичили 13 з 16 штрафів.} \\
\texttt{Інший гравець надіслав це повідомлення:} \\
\\ 
\\ 
\texttt{"зачекай, я перевірю свої координати. Т: (1,1), A: (2,1), D: (3,1), O: (1,2), P: (2,2), S: (3,2), M: (1,3). Все правильно?"} \\
\\ 
\\ 
\texttt{Яка ваша наступна команда?} \\
            }
        }
    }
    & \\ \\

    \theutterance \stepcounter{utterance}  
    & & & \multicolumn{4}{p{0.6\linewidth}}{
        \cellcolor[rgb]{0.9,0.9,0.9}{
            \makecell[{{p{\linewidth}}}]{
                \texttt{\tiny{[GM$\langle$P2]}}
                \texttt{ПЕРЕСУНЬ: O, (1, 2)} \\
            }
        }
    }
     \\ \\

    \theutterance \stepcounter{utterance}  
    & & & \multicolumn{2}{p{0.3\linewidth}}{
        \cellcolor[rgb]{0.9,0.9,0.9}{
            \makecell[{{p{\linewidth}}}]{
                \texttt{\tiny{[GM$|$GM]}}
                \texttt{ПЕРЕСУНЬ: O, (1, 2)} \\
            }
        }
    }
    & & \\ \\

    \theutterance \stepcounter{utterance}  
    & & & \multicolumn{2}{p{0.3\linewidth}}{
        \cellcolor[rgb]{0.9,0.9,0.9}{
            \makecell[{{p{\linewidth}}}]{
                \texttt{\tiny{[GM$|$GM]}}
                \texttt{O moved to (1, 2)} \\
            }
        }
    }
    & & \\ \\

    \theutterance \stepcounter{utterance}  
    & & \multicolumn{4}{p{0.6\linewidth}}{
        \cellcolor[rgb]{0.9,0.9,0.9}{
            \makecell[{{p{\linewidth}}}]{
                \texttt{\tiny{[P1$\langle$GM]}}
                \texttt{Ваше повідомлення передано іншому гравцю.} \\
\\ 
\\ 
\texttt{Зараз ви граєте раунд 17 з максимум 28.} \\
\texttt{Разом ви накопичили 13 з 16 штрафів.} \\
\texttt{Інший гравець перемістив об'єкт на своїй сітці.} \\
\\ 
\\ 
\texttt{Яка ваша наступна команда?} \\
            }
        }
    }
    & \\ \\

    \theutterance \stepcounter{utterance}  
    & \multicolumn{4}{p{0.6\linewidth}}{
        \cellcolor[rgb]{0.9,0.9,0.9}{
            \makecell[{{p{\linewidth}}}]{
                \texttt{\tiny{[P1$\rangle$GM]}}
                \texttt{СКАЖИ: зачекай, я перевірю свої координати. Т: (1,1), A: (2,1), D: (3,1), O: (1,2), P: (2,2), S: (3,2), M: (1,3). Все правильно?} \\
            }
        }
    }
    & & \\ \\

    \theutterance \stepcounter{utterance}  
    & & & \multicolumn{2}{p{0.3\linewidth}}{
        \cellcolor[rgb]{0.9,0.9,0.9}{
            \makecell[{{p{\linewidth}}}]{
                \texttt{\tiny{[GM$|$GM]}}
                \texttt{СКАЖИ: зачекай, я перевірю свої координати. Т: (1,1), A: (2,1), D: (3,1), O: (1,2), P: (2,2), S: (3,2), M: (1,3). Все правильно?} \\
            }
        }
    }
    & & \\ \\

    \theutterance \stepcounter{utterance}  
    & & \multicolumn{4}{p{0.6\linewidth}}{
        \cellcolor[rgb]{0.9,0.9,0.9}{
            \makecell[{{p{\linewidth}}}]{
                \texttt{\tiny{[GM$\rangle$P2]}}
                \texttt{Об'єкт 'O' успішно переміщено в (1,2). Ваша оновлена сітка виглядає так:} \\
\texttt{```} \\
\texttt{ 1234567} \\
\texttt{╔═══════╗} \\
\texttt{║TAD◌◌◌◌║ 1} \\
\texttt{║OPS◌◌◌◌║ 2} \\
\texttt{║◌M◌◌◌┌─╢ 3} \\
\texttt{╟──┐◌◌│◌║ 4} \\
\texttt{║◌◌├──┼─╢ 5} \\
\texttt{║◌◌│◌◌│◌║ 6} \\
\texttt{║◌◌│◌◌│◌║ 7} \\
\texttt{╚══╧══╧═╝} \\
\texttt{```} \\
\texttt{Зараз ви граєте раунд 18 з максимум 28.} \\
\texttt{Разом ви накопичили 13 з 16 штрафів.} \\
\texttt{Інший гравець надіслав це повідомлення:} \\
\\ 
\\ 
\texttt{"зачекай, я перевірю свої координати. Т: (1,1), A: (2,1), D: (3,1), O: (1,2), P: (2,2), S: (3,2), M: (1,3). Все правильно?"} \\
\\ 
\\ 
\texttt{Яка ваша наступна команда?} \\
            }
        }
    }
    & \\ \\

    \theutterance \stepcounter{utterance}  
    & & & \multicolumn{4}{p{0.6\linewidth}}{
        \cellcolor[rgb]{0.9,0.9,0.9}{
            \makecell[{{p{\linewidth}}}]{
                \texttt{\tiny{[GM$\langle$P2]}}
                \texttt{ПЕРЕСУНЬ: M, (1, 3)} \\
            }
        }
    }
     \\ \\

    \theutterance \stepcounter{utterance}  
    & & & \multicolumn{2}{p{0.3\linewidth}}{
        \cellcolor[rgb]{0.9,0.9,0.9}{
            \makecell[{{p{\linewidth}}}]{
                \texttt{\tiny{[GM$|$GM]}}
                \texttt{ПЕРЕСУНЬ: M, (1, 3)} \\
            }
        }
    }
    & & \\ \\

    \theutterance \stepcounter{utterance}  
    & & & \multicolumn{2}{p{0.3\linewidth}}{
        \cellcolor[rgb]{0.9,0.9,0.9}{
            \makecell[{{p{\linewidth}}}]{
                \texttt{\tiny{[GM$|$GM]}}
                \texttt{M moved to (1, 3)} \\
            }
        }
    }
    & & \\ \\

    \theutterance \stepcounter{utterance}  
    & & \multicolumn{4}{p{0.6\linewidth}}{
        \cellcolor[rgb]{0.9,0.9,0.9}{
            \makecell[{{p{\linewidth}}}]{
                \texttt{\tiny{[P1$\langle$GM]}}
                \texttt{Ваше повідомлення передано іншому гравцю.} \\
\\ 
\\ 
\texttt{Зараз ви граєте раунд 18 з максимум 28.} \\
\texttt{Разом ви накопичили 13 з 16 штрафів.} \\
\texttt{Інший гравець перемістив об'єкт на своїй сітці.} \\
\\ 
\\ 
\texttt{Яка ваша наступна команда?} \\
            }
        }
    }
    & \\ \\

    \theutterance \stepcounter{utterance}  
    & \multicolumn{4}{p{0.6\linewidth}}{
        \cellcolor[rgb]{0.9,0.9,0.9}{
            \makecell[{{p{\linewidth}}}]{
                \texttt{\tiny{[P1$\rangle$GM]}}
                \texttt{СКАЖИ: зачекай, я перевірю свої координати. Т: (1,1), A: (2,1), D: (3,1), O: (1,2), P: (2,2), S: (3,2), M: (1,3). Все правильно?} \\
            }
        }
    }
    & & \\ \\

    \theutterance \stepcounter{utterance}  
    & & & \multicolumn{2}{p{0.3\linewidth}}{
        \cellcolor[rgb]{0.9,0.9,0.9}{
            \makecell[{{p{\linewidth}}}]{
                \texttt{\tiny{[GM$|$GM]}}
                \texttt{СКАЖИ: зачекай, я перевірю свої координати. Т: (1,1), A: (2,1), D: (3,1), O: (1,2), P: (2,2), S: (3,2), M: (1,3). Все правильно?} \\
            }
        }
    }
    & & \\ \\

    \theutterance \stepcounter{utterance}  
    & & \multicolumn{4}{p{0.6\linewidth}}{
        \cellcolor[rgb]{0.9,0.9,0.9}{
            \makecell[{{p{\linewidth}}}]{
                \texttt{\tiny{[GM$\rangle$P2]}}
                \texttt{Об'єкт 'M' успішно переміщено в (1,3). Ваша оновлена сітка виглядає так:} \\
\texttt{```} \\
\texttt{ 1234567} \\
\texttt{╔═══════╗} \\
\texttt{║TAD◌◌◌◌║ 1} \\
\texttt{║OPS◌◌◌◌║ 2} \\
\texttt{║M◌◌◌◌┌─╢ 3} \\
\texttt{╟──┐◌◌│◌║ 4} \\
\texttt{║◌◌├──┼─╢ 5} \\
\texttt{║◌◌│◌◌│◌║ 6} \\
\texttt{║◌◌│◌◌│◌║ 7} \\
\texttt{╚══╧══╧═╝} \\
\texttt{```} \\
\texttt{Зараз ви граєте раунд 19 з максимум 28.} \\
\texttt{Разом ви накопичили 13 з 16 штрафів.} \\
\texttt{Інший гравець надіслав це повідомлення:} \\
\\ 
\\ 
\texttt{"зачекай, я перевірю свої координати. Т: (1,1), A: (2,1), D: (3,1), O: (1,2), P: (2,2), S: (3,2), M: (1,3). Все правильно?"} \\
\\ 
\\ 
\texttt{Яка ваша наступна команда?} \\
            }
        }
    }
    & \\ \\

    \theutterance \stepcounter{utterance}  
    & & & \multicolumn{4}{p{0.6\linewidth}}{
        \cellcolor[rgb]{0.9,0.9,0.9}{
            \makecell[{{p{\linewidth}}}]{
                \texttt{\tiny{[GM$\langle$P2]}}
                \texttt{СКАЖИ: завершено!} \\
            }
        }
    }
     \\ \\

    \theutterance \stepcounter{utterance}  
    & & & \multicolumn{2}{p{0.3\linewidth}}{
        \cellcolor[rgb]{0.9,0.9,0.9}{
            \makecell[{{p{\linewidth}}}]{
                \texttt{\tiny{[GM$|$GM]}}
                \texttt{СКАЖИ: завершено!} \\
            }
        }
    }
    & & \\ \\

    \theutterance \stepcounter{utterance}  
    & & & \multicolumn{2}{p{0.3\linewidth}}{
        \cellcolor[rgb]{0.9,0.9,0.9}{
            \makecell[{{p{\linewidth}}}]{
                \texttt{\tiny{[GM$|$GM]}}
                \texttt{true} \\
            }
        }
    }
    & & \\ \\

    \theutterance \stepcounter{utterance}  
    & & & \multicolumn{2}{p{0.3\linewidth}}{
        \cellcolor[rgb]{0.9,0.9,0.9}{
            \makecell[{{p{\linewidth}}}]{
                \texttt{\tiny{[GM$|$GM]}}
                \texttt{Initial states:} \\
\texttt{```} \\
\texttt{Player 1:} \\
\texttt{ 1234567} \\
\texttt{╔═══════╗} \\
\texttt{║◌T◌◌◌◌◌║ 1} \\
\texttt{║A◌◌◌◌◌D║ 2} \\
\texttt{║◌◌◌◌◌┌─╢ 3} \\
\texttt{╟──┐O◌│◌║ 4} \\
\texttt{║◌P├──┼─╢ 5} \\
\texttt{║◌S│M◌│◌║ 6} \\
\texttt{║◌◌│◌◌│◌║ 7} \\
\texttt{╚══╧══╧═╝} \\
\\ 
\\ 
\texttt{Player 2:} \\
\texttt{ 1234567} \\
\texttt{╔═══════╗} \\
\texttt{║◌◌P◌A◌O║ 1} \\
\texttt{║S◌◌◌M◌◌║ 2} \\
\texttt{║◌◌◌◌D┌─╢ 3} \\
\texttt{╟──┐◌◌│◌║ 4} \\
\texttt{║◌◌├──┼─╢ 5} \\
\texttt{║◌◌│◌◌│T║ 6} \\
\texttt{║◌◌│◌◌│◌║ 7} \\
\texttt{╚══╧══╧═╝} \\
\\ 
\texttt{```} \\
            }
        }
    }
    & & \\ \\

    \theutterance \stepcounter{utterance}  
    & & & \multicolumn{2}{p{0.3\linewidth}}{
        \cellcolor[rgb]{0.9,0.9,0.9}{
            \makecell[{{p{\linewidth}}}]{
                \texttt{\tiny{[GM$|$GM]}}
                \texttt{End states:} \\
\texttt{```} \\
\texttt{Player 1:} \\
\texttt{ 1234567} \\
\texttt{╔═══════╗} \\
\texttt{║TAD◌◌◌◌║ 1} \\
\texttt{║OPS◌◌◌◌║ 2} \\
\texttt{║M◌◌◌◌┌─╢ 3} \\
\texttt{╟──┐◌◌│◌║ 4} \\
\texttt{║◌◌├──┼─╢ 5} \\
\texttt{║◌◌│◌◌│◌║ 6} \\
\texttt{║◌◌│◌◌│◌║ 7} \\
\texttt{╚══╧══╧═╝} \\
\\ 
\\ 
\texttt{Player 2:} \\
\texttt{ 1234567} \\
\texttt{╔═══════╗} \\
\texttt{║TAD◌◌◌◌║ 1} \\
\texttt{║OPS◌◌◌◌║ 2} \\
\texttt{║M◌◌◌◌┌─╢ 3} \\
\texttt{╟──┐◌◌│◌║ 4} \\
\texttt{║◌◌├──┼─╢ 5} \\
\texttt{║◌◌│◌◌│◌║ 6} \\
\texttt{║◌◌│◌◌│◌║ 7} \\
\texttt{╚══╧══╧═╝} \\
\\ 
\texttt{```} \\
            }
        }
    }
    & & \\ \\

    \theutterance \stepcounter{utterance}  
    & & & \multicolumn{2}{p{0.3\linewidth}}{
        \cellcolor[rgb]{0.9,0.9,0.9}{
            \makecell[{{p{\linewidth}}}]{
                \texttt{\tiny{[GM$|$GM]}}
                \texttt{* success: True} \\
\texttt{* lose: False} \\
\texttt{* aborted: False} \\
\texttt{{-}{-}{-}{-}{-}{-}{-}} \\
\texttt{* Shifts: 12.00} \\
\texttt{* Max Shifts: 12.00} \\
\texttt{* Min Shifts: 6.00} \\
\texttt{* End Distance Sum: 0.00} \\
\texttt{* Init Distance Sum: 30.04} \\
\texttt{* Expected Distance Sum: 29.33} \\
\texttt{* Penalties: 13.00} \\
\texttt{* Max Penalties: 16.00} \\
\texttt{* Rounds: 19.00} \\
\texttt{* Max Rounds: 28.00} \\
\texttt{* Object Count: 7.00} \\
            }
        }
    }
    & & \\ \\

\end{supertabular}
}

\end{document}
